% SPDX-FileCopyrightText: 2026 Will Cook
% SPDX-License-Identifier: CC-BY-4.0
%
% A mathematician-facing exposition of the Erdos249257 Lean
% development. Every formalised statement links to the exact Lean source line
% at the pinned formal-source checkpoint. Build with
% `tectonic erdos249-257-main-paper.tex` or `make`.
\documentclass[11pt]{article}

\usepackage[T1]{fontenc}
\usepackage[utf8]{inputenc}
\usepackage[a4paper,margin=1in]{geometry}
\usepackage{amsmath,amssymb,amsthm}
\usepackage{booktabs}
\usepackage{enumitem}
\usepackage[dvipsnames]{xcolor}
\usepackage[colorlinks=true,linkcolor=black,urlcolor=RoyalBlue,citecolor=black]{hyperref}
\hypersetup{
  pdftitle={Tail Certificates and Achievement-Set Geometry for Erdős Problems 249 and 257},
  pdfauthor={Will Cook},
  pdfsubject={Main gateway paper for the Erdős 249/257 Lean formalisation},
}

% ---------------------------------------------------------------------------
% Drill-down machinery: each formal anchor links to the GitHub blob at the
% pinned formal-source commit. \commit names that checkpoint; scripts/check_release.py in
% the repository verifies that it matches docs/claims.json and that every
% \lref/\lrefx/\lword/\lloc target resolves to the named declaration at that line.
% ---------------------------------------------------------------------------
\newcommand{\commit}{a16aa24af5236b6508a34b8a4c4b41fdcf620a98}
\newcommand{\repobase}{https://github.com/wcook04/plectis-lean-erdos249-257/blob/\commit}
\newcommand{\sourceurl}{https://github.com/wcook04/plectis-lean-erdos249-257/tree/\commit}
\newcommand{\rootbase}{https://github.com/wcook04/plectis-lean-erdos249-257/blob/main}
\newcommand{\PK}{Erdos249257}
% SPDX-FileCopyrightText: 2026 Will Cook
% SPDX-License-Identifier: CC-BY-4.0
% Generated by scripts/build_paper_module_aliases.py; do not edit.
% Each CamelCase component contributes at most its first three characters.
\newcommand{\modulesigil}[1]{\csname modulesigil@#1\endcsname}
\expandafter\def\csname modulesigil@ActualForeignResidueProjection.lean\endcsname{ActForResPro}
\expandafter\def\csname modulesigil@AdelicHeightObstruction.lean\endcsname{AdeHeiObs}
\expandafter\def\csname modulesigil@BooleanMobiusCarry.lean\endcsname{BooMobCar}
\expandafter\def\csname modulesigil@CampbellShiftSynchronization.lean\endcsname{CamShiSyn}
\expandafter\def\csname modulesigil@CarrySurvivorExtinction.lean\endcsname{CarSurExt}
\expandafter\def\csname modulesigil@CertificateKernel.lean\endcsname{CerKer}
\expandafter\def\csname modulesigil@DiagonalFlexibleOddWindowAffine.lean\endcsname{DiaFleOddWinAff}
\expandafter\def\csname modulesigil@DiagonalFlexibleOddWindowSupply.lean\endcsname{DiaFleOddWinSup}
\expandafter\def\csname modulesigil@DiagonalFreshLossBridge.lean\endcsname{DiaFreLosBri}
\expandafter\def\csname modulesigil@DiagonalPincerCertificateT64Endpoint.lean\endcsname{DiaPinCerT64End}
\expandafter\def\csname modulesigil@DiagonalPincerCertificatesT64.lean\endcsname{DiaPinCerT64}
\expandafter\def\csname modulesigil@DiagonalPincerDecomposition.lean\endcsname{DiaPinDec}
\expandafter\def\csname modulesigil@DiagonalPincerPrimeCertificates/ClosureT64.lean\endcsname{DiaPinPriCerT64}
\expandafter\def\csname modulesigil@FreshPrimeDeficitDecomposition.lean\endcsname{FrePriDefDec}
\expandafter\def\csname modulesigil@FullTargetPrimeAdjunctionNoGo.lean\endcsname{FulTarPriAdjNoGo}
\expandafter\def\csname modulesigil@GapFareyBound.lean\endcsname{GapFarBou}
\expandafter\def\csname modulesigil@GcdMomentCalculus.lean\endcsname{GcdMomCal}
\expandafter\def\csname modulesigil@GenericTailOrbitRigidity.lean\endcsname{GenTaiOrbRig}
\expandafter\def\csname modulesigil@GreedyAchievementSet.lean\endcsname{GreAchSet}
\expandafter\def\csname modulesigil@LcmConeFlatness.lean\endcsname{LcmConFla}
\expandafter\def\csname modulesigil@LcmConeNonflat.lean\endcsname{LcmConNon}
\expandafter\def\csname modulesigil@LcmDiagonalReduction.lean\endcsname{LcmDiaRed}
\expandafter\def\csname modulesigil@MersenneLambertLadder.lean\endcsname{MerLamLad}
\expandafter\def\csname modulesigil@MersenneShadowDenominatorGrowth.lean\endcsname{MerShaDenGro}
\expandafter\def\csname modulesigil@MersenneTailAtoms.lean\endcsname{MerTaiAto}
\expandafter\def\csname modulesigil@MobiusSignSupportNoGo.lean\endcsname{MobSigSupNoGo}
\expandafter\def\csname modulesigil@PowerTwoBitLift.lean\endcsname{PowTwoBitLif}
\expandafter\def\csname modulesigil@PowerTwoCenteredBitLift.lean\endcsname{PowTwoCenBitLif}
\expandafter\def\csname modulesigil@PrimitiveDeterminantLift.lean\endcsname{PriDetLif}
\expandafter\def\csname modulesigil@PrimitiveWeightCertificate.lean\endcsname{PriWeiCer}
\expandafter\def\csname modulesigil@RationalSupportCarrySkeleton.lean\endcsname{RatSupCarSke}
\expandafter\def\csname modulesigil@ResidualGaugeObstruction.lean\endcsname{ResGauObs}
\expandafter\def\csname modulesigil@SignedQMomentObstruction.lean\endcsname{SigQMomObs}
\expandafter\def\csname modulesigil@SquareCRTCube.lean\endcsname{SquCRTCub}
\expandafter\def\csname modulesigil@SquaredMersenneDiagonalEnclosure.lean\endcsname{SquMerDiaEnc}
\expandafter\def\csname modulesigil@SternBrocotRunGeometry.lean\endcsname{SteBroRunGeo}
\expandafter\def\csname modulesigil@SublogDivisorCoverage.lean\endcsname{SubDivCov}
\expandafter\def\csname modulesigil@TotientCarryKernelRigidity.lean\endcsname{TotCarKerRig}
\expandafter\def\csname modulesigil@TotientMahlerDefect.lean\endcsname{TotMahDef}
\expandafter\def\csname modulesigil@TotientTailPeriodKiller.lean\endcsname{TotTaiPerKil}

% Formal source coordinates and exact declaration names remain in the
% hyperlink targets and verifier manifest. Default links turn the declaration
% name into a spaced, lower-case label; \lword supplies a shorter authored
% label when the surrounding exposition needs one.
\ExplSyntaxOn
\cs_new_protected:Npn \paperleanlabel #1
  {
    \tl_set:Nx \l_tmpa_tl { \tl_to_str:n {#1} }
    \regex_replace_all:nnN { _+ } { \c{space} } \l_tmpa_tl
    \regex_replace_all:nnN
      { ([a-z0-9])([A-Z]) }
      { \1\c{space}\2 }
      \l_tmpa_tl
    \tl_set:Nx \l_tmpa_tl { \tl_lower_case:n { \tl_use:N \l_tmpa_tl } }
    \tl_use:N \l_tmpa_tl
  }
\ExplSyntaxOff
\newcommand{\leanlabel}[1]{\paperleanlabel{#1}}
\newcommand{\lref}[3]{\href{\repobase/\PK/#1\#L#2}{\leanlabel{#3}}}
\newcommand{\lrefx}[3]{\href{\repobase/\PK/#1\#L#2}{\leanlabel{#3}}}
\newcommand{\lword}[4]{\href{\repobase/\PK/#1\#L#2}{#4}}
\newcommand{\lloc}[2]{\href{\repobase/\PK/#1\#L#2}{Lean source}}
% \rref{rootfile}: link a repository-root navigation file at the public main branch.
\newcommand{\rref}[1]{\href{\rootbase/\detokenize{#1}}{repository record}}
\emergencystretch=2em

\newtheorem{theorem}{Theorem}[section]
\newtheorem{proposition}[theorem]{Proposition}
\newtheorem{corollary}[theorem]{Corollary}
\theoremstyle{definition}
\newtheorem{definition}[theorem]{Definition}
\newtheorem{example}[theorem]{Example}
\newtheorem{problem}[theorem]{Problem}
\theoremstyle{remark}
\newtheorem*{remark}{Remark}

\newcommand{\N}{\mathbb{N}}
\newcommand{\Z}{\mathbb{Z}}
\newcommand{\Q}{\mathbb{Q}}
\newcommand{\R}{\mathbb{R}}
\newcommand{\ph}{\varphi}
\newcommand{\lcm}{\operatorname{lcm}}
\newcommand{\Irr}{\text{ irrational}}
\newcommand{\Qzero}{Q_0}

\title{\vspace{-4em}Tail Certificates and Achievement-Set Geometry\\
{\large for Erd\H{o}s Problems 249 and 257}}
\author{Will Cook \qquad July 2026}
\date{}

\begin{document}
\sloppy
\maketitle

\begin{abstract}
\noindent
For the totient constant
$S=\sum_{n\ge1}\ph(n)/2^n$, one certified Farey window excludes every rational
representation $S=p/q$ with
$q\le\Qzero\approx7.96\times10^{34}$.  We also prove that irrationality of
$S$ is equivalent to non-integrality of every positive binary tail difference,
with a complete finite certificate at each fixed pair of parameters; one
integral difference would instead force rationality.  The missing input is an
unbounded supply of these certificates, and the checks through $t=64$ are
finite evidence only.  For Erd\H{o}s~\#257, we formalise Erd\H{o}s's classical
full-support theorem in every integer base and several structured-support
cases.  At base $2$, the Mersenne achievement set is compact, perfect, nowhere
dense, and of measure one.  Any infinite rational-valued support must have
unbounded positive carry states and sublogarithmic divisor-coverage gaps.
Membership of $1/2$, which would refute the universal statement, is equivalent
to infinitely many greedy skips.  The upper transition and the middle
coordinate $-3$ cannot be final.  To exclude a final skip altogether, it
remains to rule out the exceptional coordinates $-2,-1$ and prove the exact
future-tail inequality for every other non-$-3$ middle transition.  Neither
problem is settled.
\end{abstract}

\section{Introduction}\label{sec:intro}

The two questions are whether
$S=\sum_{n\ge1}\ph(n)/2^n$ is irrational (\#249), and whether
$\sum_{n\in A}(2^n-1)^{-1}$ is irrational for every infinite
$A\subseteq\N$ (\#257).  The formulations occur in the Erd\H{o}s--Graham
collection and Erd\H{o}s's later survey~\cite{erdosgraham1980,erdos1988};
the numbering follows the maintained catalogue~\cite{erdosproblems}.

For \#257, proving irrationality for every infinite support is the universal
direction; constructing an infinite support of sum $1/2$ would refute it.
The results below contain unconditional positive cases and an exact conditional
counterexample route, but they reach neither endpoint.  For \#249, and for the
half-value counterexample branch of \#257, tail and carry identities separate
exact finite arithmetic from one analytic or combinatorial cofinality
statement.  The universal \#257 direction remains open for arbitrary infinite
supports.  The status classes are:

\begin{center}
\scriptsize
\begin{tabular}{@{}p{3cm}p{5.3cm}p{5.3cm}@{}}
\toprule
status & Erd\H{o}s \#249 & Erd\H{o}s \#257\\
\midrule
\textbf{\mbox{Classical}} &
Lambert-series and convolution identities used to represent $S$. &
Full support in every base $b\ge2$; pairwise-coprime and periodic support
families.\\[2pt]
\textbf{\mbox{Unconditional/exact}} &
One window gives $q>\Qzero$ for every rational $S=p/q$; irrationality is
equivalent to non-integrality of every positive tail difference and to
pointwise finite certificates. &
Measure-one nowhere-dense achievement geometry; exact rational carry
certificates; unbounded states and sublogarithmic zero gaps over infinite
rational supports; exact classification of $1/2$.\\[2pt]
\textbf{\mbox{Verified finite}} &
28 diagonal certificates through $t=64$ and further bounded exact
verifications. &
Local band and window identities, bounded reset rows, and explicit finite
counterexamples.\\[2pt]
\textbf{\mbox{Conditional}} &
A cofinal certificate or scalar-slack supply implies irrationality. &
Cofinal seam, window, strip, or terminal hypotheses imply an infinite
support of sum $1/2$.\\[2pt]
\textbf{\mbox{Open}} &
An unbounded certificate supply. &
Universal direction: all infinite supports.  Counterexample direction:
cofinal false terminals; concretely, the two final-skip obligations in
Section~\ref{sec:eb}.\\
\bottomrule
\end{tabular}
\end{center}

\paragraph{Reading map.}
Section~\ref{sec:spines} states the two exact proof spines and their unresolved
hypotheses.  Section~\ref{sec:ladder} gives the common Mersenne--Lambert
transform; Sections~\ref{sec:eb} and~\ref{sec:249} give the theorem-level
detail.  Section~\ref{sec:further} states the remaining questions.
Appendices~\ref{app:half-technical} and~\ref{sec:carry} contain the local
half-value and rational-carry machinery; the later appendices contain the
totient-certificate reductions and the independent Lambert, probability, and
finite-algebraic structure.  Linked mathematical phrases open exact Lean
proofs.  The companion develops the affine and first-harmonic \#249 theory
without changing its open status.

\medskip
\noindent\textbf{Keywords.} irrationality; Erd\H{o}s--Borwein constant;
Euler totient; Lambert series; achievement sets; binary carry systems; Lean~4.
\qquad
\textbf{MSC 2020.} 11J72 (primary); 11A25, 68V20 (secondary).

\section{The two exact spines}\label{sec:spines}

\subsection{Erd\H{o}s \#249: an exact certificate normal form}
\label{res:249story}

Write
\[
 R_N=\sum_{m\geq1}\frac{\ph(N+m)}{2^m}.
\]
The formal source proves the exact chain
\[
\begin{split}
 S\notin\Q
 &\quad\Longleftrightarrow\quad
   R_{N+h}-R_N\notin\Z
   \quad\text{for every }h>0\text{ and every }N\\
 &\quad\Longleftrightarrow\quad
   \text{for every }h>0,N\text{ there is an }L
   \text{ with }\textsf{certifiedKill}(h,N,L).
\end{split}
\]
Here \(\textsf{certifiedKill}(h,N,L)\) is a finite residue computation with a
rigorous tail allowance.  The second equivalence is completeness of the
certificate: at fixed \(h,N\), some finite depth succeeds exactly when the
tail difference is non-integral.  Conversely, integrality of one
positive-shift difference forces \(S\) to be rational.

This normal form is accompanied by an unconditional theorem: if \(S=p/q\),
then
\[
 q>\Qzero
 =79\,639\,646\,646\,701\,375\,323\,355\,774\,875\,831\,053.
\]
The exact open statement is an unbounded supply of non-integral tail
differences, equivalently of finite certificates.  Corollary~\ref{res:diag}
gives a sufficient lcm-diagonal form, and finite diagonal instances are
verified through \(t=64\), but no unbounded supply is proved.  The displayed
chain is formalised by the
\lword{LcmConeFlatness.lean}{386}{irrational_totient_series_iff_all_tail_diffs_nonintegral}{tail-difference equivalence},
the
\lword{LcmConeFlatness.lean}{399}{irrational_totient_series_iff_pointwise_certificates}{pointwise certificate equivalence},
and the fixed-parameter
\lword{LcmConeFlatness.lean}{316}{exists_certifiedKill_iff_tail_diff_notMem_int}{certificate-completeness theorem}.
The unconditional finite bound is the
\lword{CertificateKernel.lean}{18056}{tsum_totient_div_pow_two_ne_ratCast_of_den_le_79639646646701375323355774875831053}{denominator exclusion}.

\subsection{Erd\H{o}s \#257: the exact half-value counterexample spine}
\label{res:halfstory}

Let
\[
  \mathcal A=\left\{\sum_{n\in A}\frac1{2^n-1}:A\subseteq\N,\ 0\notin A\right\}.
\]
The weights are strictly superincreasing, so the normalized support of a value
is unique.  For \(1/2\), the following are equivalent:
\[
\begin{aligned}
\frac12\in\mathcal A
&\quad\Longleftrightarrow\quad
\text{the greedy skipped support is infinite}\\
&\quad\Longleftrightarrow\quad
\text{false successor terminals occur beyond every bound}.
\end{aligned}
\]
A finite normalized support cannot have value \(1/2\).  Thus membership would
produce an infinite support of rational sum and refute the universal statement
of Erd\H{o}s~\#257.  The required counterexample condition is therefore exact:
cofinally many false terminals, equivalently no last greedy skip.

More generally, rationality over an infinite support forces an unbounded
positive integer carry orbit.  Further exact consequences are
reciprocal-mass lower bounds and sublogarithmic zero-window bounds for the divisor-count
coefficient of any rational-valued support.  These are unconditional rigidity
theorems, not a proof that rational infinite supports cannot exist.

The strongest current last-skip contradiction is conditional.  At the seam
following a putative last skip, the upper transition is impossible, and among
the middle coordinates \(-3,-2,-1\), the coordinate \(-3\) is impossible.
The remaining proof must exclude the exceptional coordinates \(-2,-1\) and
establish the exact future-tail inequality at every other non-\(-3\) middle
transition.  Neither assertion is proved for the actual greedy orbit.  The
membership chain is formalised by the
\lword{GreedyAchievementSet.lean}{2514}{half_mem_mersenneAchievementSet_iff_greedySkippedSupport_infinite}{infinite-skip equivalence},
\lword{HalfCylinderHalfMembershipClassification.lean}{126}{half_mem_mersenneAchievementSet_iff_unboundedTerminalFalse}{unbounded-terminal equivalence},
\lword{HalfCylinderHalfMembershipClassification.lean}{235}{half_mem_mersenneAchievementSet_iff_no_lastHalfGreedySkip}{no-last-skip equivalence},
and the
\lword{HalfCarryReachability.lean}{589}{finite_boolSupport_ne_half}{finite-support exclusion}.
The conditional residual is the
\lword{HalfCylinderLastProducerContradiction.lean}{387}{half_mem_mersenneAchievementSet_of_middleProducerTailEscapeExceptNegThree}{remaining-tail implication}.
Its local transition exclusions are stated and linked where they are used in
Theorem~\ref{res:fatalright}; Appendix~\ref{app:half-technical} contains the
supporting coordinate analysis.

\paragraph{Prior-art boundary.}
The full-support theorem is Erd\H{o}s's, and Borwein proved a broader
irrationality theorem containing the same Mersenne--Lambert
specialisation~\cite{erdos1948,borwein1992}.  Pairwise-coprime supports are
classical~\cite{erdos1968}, and Luca--Tachiya treat periodic rational
coefficients by another method~\cite{lucatachiya}.  Kova\v{c}--Tao supply the
strict-tail achievement-set context~\cite{kovactao}; Wang--Grau Ribas supply
the closest recorded integral carry recurrence~\cite{wanggrauribas}.
Crandall and Campbell concern adjacent binary-expansion questions
~\cite{crandall,campbell}.  The explicit denominator bound, certificate
equivalence, and last-skip reduction have no exact antecedent in the
maintained comparison record; that absence is not itself a novelty claim.
The theorem-specific comparison record is maintained with the release.

\section{The Mersenne--Lambert ladder}\label{sec:ladder}

Define the Lambert-type transform of an arithmetic function $f$ by
\[
L(f)\;=\;\sum_{n\ge1}\frac{f(n)}{2^n-1}\;=\;\sum_{m\ge1}\frac{(f*\mathbf 1)(m)}{2^m},
\qquad (f*\mathbf 1)(m)=\sum_{d\mid m} f(d),
\]
the second equality being the standard Lambert rearrangement (one Dirichlet
convolution by the constant $\mathbf 1=\zeta$); Apostol supplies the classical
arithmetic-function and Dirichlet-convolution background~\cite{apostol}; Merca
introduced a modern Lambert-series factorisation theorem~\cite{merca2017}, and
Merca--Schmidt give a synthesis/unification treatment~\cite{mercaschmidt}.
Walking the ladder by convolution gives five values with
rational, irrational, open, or transcendental status.

\begin{center}
\begin{tabular}{@{}llll@{}}
\toprule
input $f$ & value & mathematical status & treatment here\\
\midrule
$\mu$ & $L(\mu)=1/2$ & rational & formalised\\
$\ph$ & $L(\ph)=2$ & rational & formalised\\
$\mathbf 1$ & $L(\mathbf 1)=E$ & irrational & formalised\\
$A=\ph*\mu$ & $L(A)=S$ & open (\#249) & identities and bounds formalised\\
$\mathrm{Id}$ & $L(\mathrm{Id})=\sum\sigma(m)/2^m$ & transcendental & cited~\cite{nesterenko}\\
\bottomrule
\end{tabular}
\end{center}

The weight $A=\ph*\mu$ is the primitive-conductor count: $A(p)=p-2$ on primes,
$A\ge0$, and $A*\zeta=\ph$. One convolution by $\zeta$ separates the open
constant $S$ from the rational $2$, and one more separates rational from
transcendental ($\ph*\zeta=\mathrm{Id}$, $\mathrm{Id}*\zeta=\sigma$). So $S$ is
the rung one convolution from the machine-checked irrational $E$: the same series
$\sum(\cdot)/(2^d-1)$, with the constant weight $1$ replaced by the nonnegative
weight $A$.

\begin{proposition}[the rational rungs]
$\sum_{d\ge1}\mu(d)/(2^d-1)=1/2$ and $\sum_{d\ge1}\ph(d)/(2^d-1)=2$, exactly.
\end{proposition}
\noindent The Lean proofs are the
\lword{MersenneLambertLadder.lean}{587}{tsum_moebius_div_two_pow_sub_one_eq_half}{M\"obius rung}
and the
\lword{MersenneLambertLadder.lean}{575}{tsum_totient_div_two_pow_sub_one_eq_two}{totient rung};
the same identities on the \#249 series shape are the
\lword{CertificateKernel.lean}{18106}{tsum_moebius_div_two_pow_sub_one_eq_half}{M\"obius restatement}
and
\lword{CertificateKernel.lean}{18101}{tsum_totient_div_two_pow_sub_one_eq_two}{totient restatement}.

\begin{proposition}[the positive lift, $L(A)=S$]\label{res:lift}
With the primitive-conductor weight $A=\ph*\mu$,
\[
\sum_{d\ge1}\frac{A(d)}{2^d-1}\;=\;\sum_{n\ge1}\frac{\ph(n)}{2^n}\;=\;S.
\]
\end{proposition}
\noindent The
\lword{CertificateKernel.lean}{18117}{tsum_primWeight_div_two_pow_sub_one_eq_totient_series}{positive-lift identity}
is formalised with the
\lword{MersenneLambertLadder.lean}{252}{primWeight}{primitive-conductor weight}.
This places $S$ in the same $\sum(\cdot)/(2^d-1)$ family as $E$, with a
nonnegative weight.

\paragraph{Probabilistic reading.}
If \(X,Y\) are independent fair-coin waiting times,
\(\Pr(X=n)=2^{-n}\), then
\[
 S=\tfrac12+\Pr(\gcd(X,Y)=1).
\]
Thus Erd\H{o}s~\#249 also asks whether the probability that two independent
fair-coin waiting times are coprime is irrational.  The exact identity is
Proposition~\ref{res:coprime}; its M\"obius-square form and the periodic and
finite-algebraic variants of the ladder are collected in
Section~\ref{sec:additional-lambert}.

\paragraph{Section consequence and residual.}
The exact contribution of this section is the coordinate change
\(L(A)=S\), together with the two rational rungs above it in the same
Lambert family.  These identities explain why the Mersenne--Lambert
coordinates are natural for both problems, but they do not imply that \(S\)
is irrational.  For \#249 the remaining issue is still the unbounded
certificate supply isolated in Section~\ref{sec:249}; for \#257 the ladder
supplies context rather than either the universal theorem or a rational
counterexample.

\section{Erd\H{o}s--Borwein-type irrationality (the \#257 direction)}\label{sec:eb}

There are two logically distinct directions.  The universal direction asks
for irrationality for every infinite support; the counterexample direction
asks for one infinite support with rational value, and the target \(1/2\)
gives an exact formulation of that problem.  This section first states the
classical and structured irrationality theorems, then gives the base-\(2\)
achievement-set geometry and the constraints forced by a hypothetical
rational value.  It ends with the exact classification of the target \(1/2\)
and the exact residual condition left by the current final-skip argument.
Local rewind, window, cylinder, and finite-stage results are
collected in Appendix~\ref{app:half-technical}.

\subsection{Full support: the Erd\H{o}s--Borwein constant, every base}

\begin{theorem}[full-support irrationality]\label{res:full}
For every integer $b\ge2$, the series $\displaystyle\sum_{n\ge1}\frac{1}{b^n-1}$ is
irrational. In particular the Erd\H{o}s--Borwein constant $E=\sum_{n\ge1}1/(2^n-1)$
is irrational.
\end{theorem}
\noindent The
\lword{CertificateKernel.lean}{8000}{irrational_erdosSum_full_support}{all-base theorem}
and its
\lword{CertificateKernel.lean}{8007}{irrational_erdosBorwein_series}{base-\(2\) corollary}
are formalised.  This is Erd\H{o}s's 1948 theorem~\cite{erdos1948}. In the ladder of
Section~\ref{sec:ladder} it is the rung $L(\mathbf 1)$, and it is the fixed
machine-checked point one convolution away from the open $S$.

The proof begins with the Lambert identity
\[
 \sum_{n\geq 1}\frac{1}{b^n-1}
   =\sum_{m\geq 1}\frac{\tau(m)}{b^m},
\]
where $\tau(m)$ is the number of positive divisors of $m$; see the
\lword{CertificateKernel.lean}{5865}{erdosSum_full_support_eq_tsum_divisor_count}{Lambert identity}.
If the right-hand side were rational with reduced denominator $q$, then a
non-integral integer multiple could not lie at distance less than $1/q$ from an
integer.  It is therefore enough, for every $q$, to construct $N$ such that
$b^N\sum_m\tau(m)b^{-m}$ is non-integral but lies within $1/q$ of an integer.
This is the formal
\lword{CertificateKernel.lean}{5792}{irrational_of_int_mul_near_int}{near-integer criterion}.

The construction of $N$ is divided into three finite parts.  First, a bounded
Bertrand argument supplies disjoint prime blocks, and the Chinese remainder
theorem chooses an arithmetic progression on which
\[
        b^r\mid \tau(N+r) \qquad (1\leq r\leq K).
\]
These divisibilities make the first $K$ terms of the shifted tail integral.
The multiplicativity step that converts prescribed prime valuations into the
displayed divisibility is the
\lword{CertificateKernel.lean}{6177}{bpow_dvd_card_divisors_of_exact_prime_block}{prime-block divisibility lemma}.
Second, divisor pairing bounds the average of the weighted middle window
\[
   \sum_{K<r\leq L}\tau(N+r)b^{L-r}
\]
along that progression.  A pigeonhole choice then selects one translate with a
small middle contribution; see the
\lword{CertificateKernel.lean}{7246}{exists_first_block_progression_with_selected_weighted_middle}{weighted-middle selection lemma}.
Third, the elementary estimate $\tau(n)\leq n$ controls the infinite tail after
$L$.  The parameters are chosen so that the middle and far-tail bounds together
are smaller than $1/q$; the shifted tail is strictly positive because every
divisor count is positive, so the resulting dilation is not an integer.

Thus the analytic series is consumed only at the Lambert identity and the
geometric tail estimate.  The remaining proof is a finite CRT construction,
divisor counting, and explicit parameter arithmetic.  The final certificate
existence theorem holds for every $b\geq2$ and every precision $q$; see the
\lword{CertificateKernel.lean}{7859}{exists_weighted_block_certificates}{uniform certificate theorem}.
Its
composition with the near-integer criterion is the theorem cited above.  This
is a formalisation of the classical full-support theorem, not a new
irrationality result.

\subsection{Named infinite-support cases}

Beyond full support, the development proves irrationality for a family of infinite
supports $A$, at every base $b\ge2$. Each is a formalised case of the \#257
statement family. None is the universal statement, and none is claimed as new
mathematics.

\begin{theorem}[support-class family]\label{res:support}
For every integer $b\ge2$, the series $\sum_{n\in A} 1/(b^n-1)$ is irrational for
each of the following supports $A$:
\begin{enumerate}[label=\textup{(\alph*)},leftmargin=2em]
\item factorial support $A=\{n! : n\ge1\}$, giving $\sum 1/(b^{n!}-1)$;
\item power-of-two support $A=\{2^n : n\ge0\}$, giving $\sum 1/(b^{2^n}-1)$;
\item multiples $A=d\N$;
\item pairwise-coprime supports with summable reciprocals (Erd\H{o}s's
condition~\cite{erdos1968});
\item eventually-periodic supports, residue classes, and the odd numbers.
\end{enumerate}
\end{theorem}
\noindent These sort into three groups. The classical case is (d): infinite
pairwise-coprime $A$ with $\sum_{a\in A}1/a<\infty$, formalised from Erd\H{o}s's
condition. The lcm-gap named cases are (a) and (b), whose support gaps outgrow the
running $\lcm$; the structured families are (c) and (e).\footnote{Luca--Tachiya
prove the broader eventual-periodic rational-coefficient theorem; its indicator
specialisation directly contains case (e)~\cite{lucatachiya}. Their
Chowla--Erd\H{o}s-style large-modulus proof is distinct from the periodic-divisor
certificate route formalised here, and the full mixed-sign coefficient theorem is
not claimed.}
\noindent Formalised case by case by the
\lword{CertificateKernel.lean}{5707}{irrational_erdosSum_factorial_support}{factorial-support theorem},
\lword{CertificateKernel.lean}{5731}{irrational_erdosSum_two_pow_support}{power-of-two-support theorem},
\lword{CertificateKernel.lean}{8775}{irrational_erdosSupportSeries_multiples}{multiple-support theorem},
\lword{CertificateKernel.lean}{10448}{irrational_erdosSupportSeries_pairwise_coprime}{pairwise-coprime-support theorem},
and
\lword{CertificateKernel.lean}{11276}{irrational_erdosSupportSeries_eventuallyPeriodic}{eventually-periodic-support theorem}.
The residue-class and odd supports in part~\textup{(e)} are direct
specialisations of the last theorem; the first two theorems already hold for
every base \(b\geq2\).

The prime support is now known unconditionally: Tao and Ter\"av\"ainen proved
that
\[
  \sum_p\frac1{2^p-1}=\sum_{n\ge1}\frac{\omega(n)}{2^n}
\]
is irrational~\cite{taoteravainen}. Their proof uses quantitative two-point
correlations of multiplicative functions and is not formalised here. This
settles an important special case, while the universal support statement
remains open.

The factorial and power-of-two cases are worth a word because they are not
reachable by a Liouville shortcut: the power-of-two gap $2^k-\lcm$ stays within a
bounded ratio of the $\lcm$, so the terms do not decay fast enough for a
transcendence-measure argument, and a uniform criterion is needed. The mechanism
is \emph{period noncollapse}: a
prime-valuation-deficit witness forces the multiplicative order of the base modulo
the relevant modulus not to collapse, which rules out the long repeated digit
blocks a rational value would require.

\subsection{The base-$2$ achievement set}

For \(n\geq1\), put
\[
 w_n=\frac1{2^n-1},
 \qquad
 T_n=\sum_{m>n}w_m,
 \qquad
 \mathcal A=\left\{\sum_{n\in A}w_n:A\subseteq\N,\ 0\notin A\right\}.
\]
For \(x\geq0\), define the greedy residuals by \(r_0(x)=x\) and
\[
 \varepsilon_n(x)=
 \begin{cases}
  1,&w_n\leq r_{n-1}(x),\\
  0,&w_n>r_{n-1}(x),
 \end{cases}
 \qquad
 r_n(x)=r_{n-1}(x)-\varepsilon_n(x)w_n.
\]

\begin{theorem}[achievement-set geometry and greedy coding]\label{res:greedy}
The following statements hold.
\begin{enumerate}[label=\textup{(\roman*)},leftmargin=2.2em]
\item \(T_n<w_n\) for every \(n\geq1\).  Hence every \(x\in\mathcal A\)
has a unique normalized support.
\item A real number \(x\) lies in \(\mathcal A\) if and only if \(x\geq0\)
and \(r_n(x)\leq T_n\) for every \(n\geq0\).  In that case
\(x=\sum_{n\geq1}\varepsilon_n(x)w_n\).
\item The set \(\mathcal A\) is compact, perfect, totally disconnected, and
nowhere dense, and its Lebesgue measure is \(1\).
\item Exact rational evaluation of the greedy recurrence is sound.  In
particular, if an exact residual exceeds an exact upper bound for the
remaining tail at some finite level, then \(x\notin\mathcal A\).
\end{enumerate}
\end{theorem}
\noindent The strict-tail inequality and greedy criterion are the
\lword{GreedyAchievementSet.lean}{180}{mersenneTail_lt_weight}{strict-tail theorem}
and
\lword{GreedyAchievementSet.lean}{1445}{mem_mersenneAchievementSet_iff_greedy_survival}{greedy membership theorem}.
Uniqueness is the
\lword{GreedyAchievementSet.lean}{1527}{positiveMersenneSupportValue_injective_normalized}{support-coding theorem};
the four topological and measure properties in part~\textup{(iii)} are proved
and linked together in Proposition~\ref{res:greedytopology}.
Kova\v{c}--Tao provide the strict-tail Cantor-set context~\cite{kovactao};
no novelty claim is made for that geometry.  Exact rational death
certificates are given in Appendix~\ref{sec:carry}.

\paragraph{Consequence for the half-value branch.}
Parts~\textup{(i)}, \textup{(ii)}, and~\textup{(iv)} make the greedy code a
canonical decision procedure, while compactness in part~\textup{(iii)} allows
cofinal finite approximations to pass to a point of \(\mathcal A\).  The
measure-one statement is global: it does not decide whether the particular
point \(1/2\) belongs to \(\mathcal A\).  That pointwise question is the
separate reduction below.

\subsection{Rigidity of a hypothetical rational support}

For \(A\subseteq\N\setminus\{0\}\), define
\[
 f_A(m)=\#\{a\in A:a\mid m\},
 \qquad
 X_A=\sum_{a\in A}\frac1{2^a-1}
     =\sum_{m\geq1}\frac{f_A(m)}{2^m},
 \qquad
 T_f(N)=\sum_{j\geq1}\frac{f(N+j)}{2^j}.
\]

\begin{theorem}[rational-support rigidity]\label{res:rigidity-main}
Suppose that \(A\neq\varnothing\) and
\[
 X_A=\frac{p}{2^c v},
 \qquad p\in\Z,\quad c\in\N,\quad v\geq1\ \text{odd}.
\]
Then \(u_N=vT_{f_A}(c+N)\) is a positive integer for every \(N\), and
\[
 u_{N+1}+v f_A(c+N+1)=2u_N. \tag{4.1}
\]
Moreover:
\begin{enumerate}[label=\textup{(\roman*)},leftmargin=2.2em]
\item if \(A\) is infinite, then \((u_N)\) is unbounded;
\item for every \(\varepsilon>0\) there is
\(B=B(\varepsilon,c,v)\) such that every interval
\(\{c+N+1,\ldots,c+N+h\}\) on which \(f_A\) vanishes satisfies
\[
 h\leq\varepsilon\log_2(N+1)+B;
\]
\item if \(\rho(A)=\sum_{a\in A}a^{-1}<\infty\), \(v>1\), and
\(h=\operatorname{ord}_v(2)\), then
\[
 \rho(A)=\frac{w}{h}
   +\lim_{M\to\infty}\frac1M\sum_{N<M}e_N\geq\frac{w}{h},
\]
where \(w\) is the number of wraps in one doubling-residue cycle and
\(u_N=ve_N+(p2^N\bmod v)\).  If \((p,v)=1\), then \(w\geq1\).
For \(v=1\) and \(A\) infinite, either \(\rho(A)\) diverges or
\(\rho(A)>1\).
\end{enumerate}
\end{theorem}
\noindent Equation~(4.1) is the exact carry recurrence.  Part (i) rules out
every bounded-state model for an infinite rational support; part (ii)
prevents logarithmically long gaps in divisor coverage; part (iii) relates
the odd denominator to the reciprocal density of the support.  These are
necessary consequences of rationality, not a contradiction.  The exact
certificate equivalence, wrap identities, and proofs are collected in
Appendix~\ref{sec:carry}.

\subsection{The target \(1/2\)}

Let
\[
 \mathcal S(x)=\{n\geq1:\varepsilon_n(x)=0\}
\]
be the set of exponents omitted by the greedy expansion.

\begin{definition}[finite greedy seam and terminal bit]\label{def:terminal-bit}
For an integer \(s\geq6\), set
\[
 q_{s,d}=\left\lfloor\frac{4^s}{2^d-1}\right\rfloor
 \quad(2\leq d<s),
 \qquad
 K_s=2^{2s-1}-2^s.
\]
Starting from \(R_{s,2}=K_s\), define successively, for \(2\leq d<s\),
\[
 \beta_{s,d}=
 \begin{cases}
  1,&q_{s,d}\leq R_{s,d},\\
  0,&q_{s,d}>R_{s,d},
 \end{cases}
 \qquad
 R_{s,d+1}=R_{s,d}-\beta_{s,d}q_{s,d}.
\]
The Boolean word \((\beta_{s,2},\ldots,\beta_{s,s-1})\) is the
\emph{finite greedy seam at row \(s\)}, and
\(\tau_s=\beta_{s,s-1}\) is its terminal bit.
\end{definition}

\noindent This is an exact integer coordinate model of the real greedy
process; Appendix~\ref{app:half-technical} develops its transition formulas
and the comparison with the real residuals.

\begin{theorem}[exact half-membership classification]\label{res:halfmembership}
The following are equivalent:
\[
\begin{aligned}
\textup{(i)}\;&\quad \frac12\in\mathcal A;\\
\textup{(ii)}\;&\quad \mathcal S(1/2)\ \text{is infinite};\\
\textup{(iii)}\;&\quad \mathcal S(1/2)\ \text{has no largest element};\\
\textup{(iv)}\;&\quad
  \text{for every }N\text{ there is }s\geq\max\{N,6\}
  \text{ with }\tau_s=0.
\end{aligned}
\]
No finite normalized support has value \(1/2\).  Consequently, any one of
\textup{(i)}--\textup{(iv)} produces an infinite support \(A\) with
\(X_A=1/2\), and therefore refutes the universal form of
Erd\H{o}s~\#257.
\end{theorem}
\noindent The greedy-support equivalence is the
\lword{GreedyAchievementSet.lean}{2514}{half_mem_mersenneAchievementSet_iff_greedySkippedSupport_infinite}{infinite-skip theorem};
the terminal forms are the
\lword{HalfCylinderHalfMembershipClassification.lean}{126}{half_mem_mersenneAchievementSet_iff_unboundedTerminalFalse}{unbounded-terminal theorem}
and
\lword{HalfCylinderHalfMembershipClassification.lean}{235}{half_mem_mersenneAchievementSet_iff_no_lastHalfGreedySkip}{no-final-skip theorem}.
Finite supports are excluded by the
\lword{HalfCarryReachability.lean}{589}{finite_boolSupport_ne_half}{finite-support theorem}.
A tempting signed-to-Boolean shortcut also fails: selecting precisely the
negative M\"obius indices overshoots \(1/2\) by at least \(1/63\).
Proposition~\ref{res:mobiussignnogo} gives the exact sign-separation identity.

\paragraph{Section consequence.}
Theorem~\ref{res:halfmembership} replaces a search over arbitrary supports by
one canonical quantifier question: does the greedy orbit of \(1/2\) skip
cofinally?  It does not answer that question.  The next subsection assumes a
possible final skip and identifies exactly what would be required to exclude
it.

\subsection{The final-skip reduction}

For each \(s\geq6\), let
\[
 G_s=\{d:2\leq d<s,\ \beta_{s,d}=1\}.
\]
Let \(H_s\) be the unique subset of \(\{2,\ldots,s-1\}\) for which
\(\sum_{d\in H_s}q_{s,d}\) is least among the subset sums strictly exceeding
\(K_s\).  The transition from row \(s\) to row \(s+1\) is called
\emph{upper}, \emph{middle}, or \emph{right} according as
\[
 G_{s+1}=H_s,\qquad G_{s+1}=G_s,\qquad
 G_{s+1}=G_s\cup\{s\},
\]
respectively.  Thus a final skipped exponent \(D\) is an upper or middle
transition followed only by right transitions.

At a middle transition \(D\geq13\), put
\[
 A_D=G_D\cup\{D\},\qquad
 P_D=2f_{G_D}(2D+1)+f_{G_D}(2D+2),
\]
and write
\[
 C_D=4R_{D,D}-P_D-4\in\Z,\qquad
 \Theta_D=\sum_{j\geq1}\frac{f_{A_D}(2D+2+j)}{2^j}.
\]
Here \(C_D\) is the affine carry at the transition and \(\Theta_D\) is its
complete future divisor-incidence tail.

For a finite set \(u\subseteq\{1,\ldots,d\}\), write
\(V(u)=\sum_{n\in u}w_n\).  Call \((u,d)\) a \emph{fatal half-gap} if
\[
 V(u)+T_{d+1}<\frac12<V(u)+w_{d+1}.
\]
Such a pair certifies that no support agreeing with \(u\) through rank \(d\)
can represent \(1/2\).

\begin{theorem}[final-skip classification and residual inequality]
\label{res:fatalright}
Nonmembership of \(1/2\) in \(\mathcal A\) is equivalent to each of the
following:
\begin{enumerate}[label=\textup{(\roman*)},leftmargin=2.2em]
\item the integer seam recurrence is eventually right;
\item a fatal half-gap exists;
\item the real greedy expansion of \(1/2\) has a final skipped exponent.
\end{enumerate}
At such a final exponent \(D\geq13\), the upper successor is impossible.
On the middle branch the value \(C_D=-3\) is impossible.  The only unresolved
exceptional cells for the current two-sided dyadic induction are
\(C_D=-2\) and \(C_D=-1\).

If every middle transition at a row \(D\geq13\) with \(C_D\neq-3\) satisfies
\[
 \Theta_D<C_D, \tag{4.2}
\]
then a final skipped exponent cannot exist, and hence \(1/2\in\mathcal A\).
\end{theorem}
\noindent The three equivalent nonmembership criteria are formalised in
\lword{HalfCylinderFatalGapRightTail.lean}{781}{seamGreedyEventuallyRight_iff_existsFatalHalfGap}{fatal-gap classification},
\lword{HalfCylinderFatalGapRightTail.lean}{787}{seamGreedyEventuallyRight_iff_half_not_mem}{nonmembership classification},
and
\lword{HalfCylinderFatalGapRightTail.lean}{802}{seamGreedyEventuallyRight_iff_exists_isLastHalfGreedySkip}{final-skip classification}.
The upper and \(C_D=-3\) exclusions are the
\lword{HalfCylinderLastProducerContradiction.lean}{250}{upperProducer_not_last}{upper-branch exclusion}
and
\lword{HalfCylinderLastProducerContradiction.lean}{315}{middleProducer_neg_three_not_last}{middle-\(-3\) exclusion}.
The three exceptional dyadic cells are identified by the
\lword{HalfCylinderMiddleCarryLowerBound.lean}{3090}{seamMiddleCoordinate_not_safe_iff_three_cells}{three-cell classification}.
The implication from~(4.2) is the
\lword{HalfCylinderLastProducerContradiction.lean}{387}{half_mem_mersenneAchievementSet_of_middleProducerTailEscapeExceptNegThree}{tail-dominance theorem}.
Since \(\Theta_D\geq0\), condition~(4.2) would in particular exclude the
remaining cells \(-2\) and \(-1\).  No proof of~(4.2) for the actual greedy
orbit is known here.

\long\def\HalfValueTechnicalAppendix{

\subsection{A band formula for the final skip between two takes}
\label{res:twothirdsband}

One local configuration admits a simple exact description in full generality.
Suppose a take at rank $b$ is followed by the next take at $c>b+1$.  Put
$q=2^b-1$, $m=2^{c-1}$, and write the residual immediately before the first
take as $1/R$.  When $0<R<q$, the final intervening skip is dyadically unsafe
if and only if
\[
  \frac{q(m-1)}{q+m-1}<R<\frac{qm}{q+m}.
\]
The interval has exact width
\[
  \frac{q^2}{(q+m)(q+m-1)}.
\]
Formalised:
\lword{HalfGreedyTwoThirdsBand.lean}{88}{postTakeUnsafeAt_iff_band}{the general unsafe-band equivalence}
and
\lword{HalfGreedyTwoThirdsBand.lean}{104}{band_width_general}{the general band-width formula}.

For a single intervening skip, $c=b+2$ and $m=2q+2$.  The band becomes
\[
  \frac{q(2q+1)}{3q+1}<R<\frac{2q(q+1)}{3q+2},
\]
with width $q^2/((3q+1)(3q+2))<1/9$, entirely between $2q/3$ and
$2q/3+2/9$.  No integral reciprocal lies in this window.  In the
corresponding positive odd integer formulation, with residual $p/(2D)$,
membership in the cleared band forces $p\ge7$.  Formalised:
\lword{HalfGreedyTwoThirdsBand.lean}{127}{singletonUnsafe_iff_twoThirdsBand}{the one-skip band equivalence},
\lword{HalfGreedyTwoThirdsBand.lean}{154}{band_width}{the one-skip width formula},
\lword{HalfGreedyTwoThirdsBand.lean}{163}{band_width_lt_ninth}{the width bound},
\lword{HalfGreedyTwoThirdsBand.lean}{185}{not_twoThirdsBand_of_int}{the integral-reciprocal exclusion},
and
\lword{HalfGreedyTwoThirdsBand.lean}{231}{seven_le_of_intBand_odd}{the odd-coordinate lower bound}.

These are localizers, not orbit-level eliminations.  They do not show that
the greedy orbit of $1/2$ avoids any such band, and they do not identify the
two remaining exceptional transition coordinates $-2,-1$ with dyadically
unsafe skips.

\subsection{Conditional criteria for structured supports}

\subsubsection{A conditional pairwise-coprime dilation criterion}\label{res:sunflower}

Let $A$ be a finite-core orthogonal-petal bouquet: outside a finite frame, its
elements have the form $c_i p_i$, where the cores $c_i$ divide one fixed
integer and the nontrivial petals $p_i$ are pairwise coprime to that frame and
to one another, with summable reciprocal mass.  The formal construction also
contains an explicit infinite example with alternating cores $2$ and $3$ and
squared prime petals; every member of this support has exactly three prime
factors with multiplicity.

For such a bouquet, the remaining input is an exact tail-selection predicate:
at every positive binary scale it chooses a carry slot with a normalized tail
at most $16$.  Lean proves that this predicate supplies the needed forced-carry
certificates and hence irrationality of the base-$2$ support series.  The
result is
\lword{SupportSunflowerDichotomy.lean}{540}{irrational_erdosSupportSeries_of_orthogonalPetalBouquet}{the conditional bouquet irrationality theorem},
with the carry bridge at
\lword{SupportSunflowerDichotomy.lean}{531}{sunflowerForcedCarrySupply_of_orthogonalPetalBouquet}{the forced-carry implication}.
The tail-selection predicate itself is not proved here.  Thus this is a
conditional reduction for one structured support family, not a further
unconditional case of Erd\H{o}s~\#257.

\subsubsection{Composite dilation defects}\label{res:compositedefect}

For a support divisor count $f_A(x)=\#\{d\in A:d\mid x\}$, multiplication by
a composite support element $a$ has an exact correction term.  Besides the
distinguished divisor $a$, the new divisors are counted by a finite defect
which excludes divisors already present at $x$.  The formal identity is
\lword{CompositeDilationDefect.lean}{30}{supportCoeff_mul_eq_add_defect}{the composite-dilation identity}.
For a finite-core orthogonal-petal bouquet, the defect at a ray is bounded by
the finite exceptional frame plus the divisor count of the petal support:
\lword{CompositeDilationDefect.lean}{151}{compositeDilationDefect_le_exceptional_add_petalCoeff}{the bouquet defect bound}.
This records the local correction needed when dilation is not prime.  It does
not prove a correlation estimate, the sunflower tail selector, or an
irrationality statement.

\subsubsection{Mixed prime-power layers}\label{res:mixedlayer}

The one-prime-power layer has a two-signature extension.  Exact $p$-adic and
$q$-adic layer operators commute, and, for distinct primes with the residual
argument coprime to $pq$, their composite extracts the iterated exact-valuation
pullback of the support coefficient.  This is formalised by
\lword{MaximalOmegaLayer.lean}{39}{mixedPrimePowerLayerTwo_supportCoeffInt}{the two-prime layer identity}.
The singleton support $\{12\}$ gives a checked $(2^2,3)$ example with value
one at residual argument $1$:
\lword{MaximalOmegaLayer.lean}{65}{mixedPrimePowerLayerTwo_twelve_fixture}{the singleton-support calculation}.
These are finite coefficient identities.  They do not provide a decimation
transport, a bounded-$\Omega$ conclusion, or an irrationality theorem.

\subsubsection{Half-trapping return carries}\label{res:halftrapping}

If rational-linear channels vanish on every relation killed by one scalar
evaluation, their evaluation matrix has rank at most one; every nontrivial
square minor therefore vanishes.  The all-rank determinant statement is
\lword{HalfTrappingReturnCarry.lean}{42}{relationInvariantLinearChannels_det_eq_zero}{the rank-one determinant criterion}.
For two reverse carries with a common coefficient word, a $1/0$ seam, and a
common following overlap, the terminal difference is a power of two times an
odd integer.  A shared nonnegative Archimedean bound then gives the matching
dyadic spacing inequality, formalised at
\lword{HalfTrappingReturnCarry.lean}{191}{overlappingReverseCarryWords_twoPow_le_realBound}{the dyadic spacing inequality}.
The argument does not produce reverse-carry presentations or an earlier
overlapping return; it supplies only the algebraic core of that conditional
criterion.

\subsubsection{Half-carry reachability}\label{res:halfcarryreachability}

There is also a finite-prefix criterion for the target $1/2$.  At an unresolved
even seam, the scalar transition reaches every integer terminal carry except
the single value $2\delta-c$; this is the exact equivalence
\lword{HalfCarryReachability.lean}{235}{evenSeamReachable_iff}{the even-seam reachability equivalence}.  A canonical
description of such a seam at every sufficiently large even depth is an
unproved hypothesis.  Under that hypothesis, the finite admissible words are
cofinal, by
\lword{HalfCarryReachability.lean}{999}{cofinalAdmissibility_of_canonicalEvenSeamSupply}{the cofinal-admissibility implication},
and compactness then produces an infinite support with sum exactly $1/2$:
\lword{HalfCarryReachability.lean}{1027}{exists_infinite_support_half_of_canonicalEvenSeamSupply}{the half-value compactness conclusion}.
Thus the formal result is a conditional reduction.  It does not establish the
canonical seam supply and does not prove Erd\H{o}s~\#257.

\paragraph{Cofinal strip returns.}  The all-level strip hypothesis can be
weakened for the canonical greedy half carry.  Its uncentred, binary-scaled
form is antitone.  Consequently, if it returns to the square-root strip
cofinally, then one sufficiently late return bounds every later scaled carry;
the strip envelope divided by $2^N$ tends to zero.  The resulting tempered
centred carry gives an infinite greedy support with support-series value
$1/2$.

\begin{proposition}[cofinal strip-return criterion]\label{res:cofinalstripreturn}
If the canonical greedy half carry returns cofinally to the stated
square-root strip, then $1/2$ belongs to the Mersenne achievement set.
\end{proposition}
\noindent Formalised:
\lword{CofinalStripReturn.lean}{84}{greedy_integerHalfCarry_scaled_tendsto_zero_of_cofinalStripReturn}{scaled-carry decay}
and
\lword{CofinalStripReturn.lean}{157}{half_mem_mersenneAchievementSet_of_cofinalStripReturn}{the strip-return membership criterion}.
The cofinal return condition itself is not established.  It is weaker than a
uniform strip bound and still does not give an unconditional membership
theorem.

\paragraph{Terminal-only cofinal approximation.}  Coherence can be removed
entirely from a further version of this reduction.  At depth $M$, take only a
normalised finite word whose terminal half carry lies in the square-root
strip.  Its finite support then has support-series value within
\[
 \frac{4M+12}{2^M}
\]
of $1/2$.  If such terminal witnesses occur cofinally, these finite values
converge to $1/2$; closedness of the Mersenne achievement set supplies the
limit.  No control of the earlier rows, and no compatibility between the
witnesses, is used.

\begin{proposition}[terminal-only cofinal criterion]\label{res:terminalonlycofinal}
Cofinal normalised terminal witnesses in the square-root strip imply that
$1/2$ belongs to the Mersenne achievement set, and hence that an infinite
support has support-series value $1/2$.
\end{proposition}
\noindent Formalised:
\lword{TerminalOnlyCofinal.lean}{68}{dist_half_erdosSupportSeries_wordSupport_le_of_terminalOnlyWitness}{the terminal approximation bound},
\lword{TerminalOnlyCofinal.lean}{134}{half_mem_mersenneAchievementSet_of_cofinalTerminalOnlyStrip}{the terminal-strip compactness criterion},
and
\lword{TerminalOnlyCofinal.lean}{192}{exists_infinite_support_half_of_cofinalTerminalOnlyStrip}{the infinite-support conclusion}.
The cofinal terminal-witness supply is not proved.  In particular, the
proposition does not establish membership of $1/2$ unconditionally.

\paragraph{Scaled terminal vanishing.}  The square-root strip is stronger
than the analytic condition used by the terminal-only argument.  Let the
depths of normalised finite words tend to infinity.  If their terminal carries,
after division by the corresponding power of $2$, tend to zero, then the
universal coefficient-tail error also tends to zero.  The finite support
values therefore converge to $1/2$, and closedness of the achievement set
again supplies the limit.  The earlier terminal-only strip condition implies
this weaker sequential hypothesis.

\begin{proposition}[scaled terminal-vanishing criterion]\label{res:terminalscaledvanishing}
A normalised terminal-word sequence with depths tending to infinity and
vanishing binary-scaled terminal carries yields an infinite support with
support-series value $1/2$.
\end{proposition}
\noindent Formalised:
\lword{TerminalOnlyScaledVanishing.lean}{24}{dist_half_erdosSupportSeries_wordSupport_le_carry_add_tail}{the scaled terminal error bound},
\lword{TerminalOnlyScaledVanishing.lean}{165}{half_mem_mersenneAchievementSet_of_terminalScaledVanishing}{the scaled-vanishing compactness criterion},
and
\lword{TerminalOnlyScaledVanishing.lean}{221}{exists_infinite_support_half_of_terminalScaledVanishing}{the scaled-vanishing support conclusion}.
No such scaled-vanishing sequence is constructed here; the result is a
conditional reduction, not an unconditional membership theorem.

\subsection{Finite-window reductions and compactness criteria}

\subsubsection{Finite shadows of the half cylinder}\label{res:halfcylindershadow}

For an admissible finite word, the residual from $1/2$ is given exactly by its
terminal half carry minus the future coefficient tail, scaled by a power of
two; see
\lword{HalfCylinderFiniteShadow.lean}{55}{halfStripAdmissible_residual_eq}{the finite-shadow residual identity}.
At a first seam, avoiding the integer strip is equivalent to full scalar
reachability of that strip:
\lword{HalfCylinderFiniteShadow.lean}{807}{evenSeam_fullIntegerStrip_iff_escape}{the seam-escape equivalence}.
The rank-three greedy example is raw dyadically safe while its seam has an
in-strip hole, so raw safety alone does not imply this escape.  Finally, an
explicit frozen-margin hypothesis implies $1/2$ belongs to the Mersenne
achievement set, by
\lword{HalfCylinderFiniteShadow.lean}{1256}{half_mem_mersenneAchievementSet_of_governedFrozenMarginProducer}{frozen-margin implication}.
That hypothesis is unproved; this is a conditional reduction, not
a proof of Erd\H{o}s~\#257.

\subsubsection{Fixed-coefficient rewind}\label{res:fixedcoeffrewind}

For a finite list of common coefficients, the canonical inverse-parent map has
an exact closed form with a power-of-two denominator:
\lword{HalfCarryCeilParentContraction.lean}{49}{rewind_eq_closedForm}{the rewind closed form}.
An interval of width at most that denominator contracts either to one ancestor
or to two adjacent ancestors.  The singleton alternative is equivalent to the
explicit dyadic phase inequality, while the other alternative is a seam pair;
see
\lword{HalfCarryCeilParentContraction.lean}{180}{rewind_singleton_or_isSeamPair}{the singleton-or-seam dichotomy}.
Thus width alone does not justify a singleton conclusion.  This is a finite
rewind lemma, and supplies neither a seam exclusion nor an Erd\H{o}s~\#257
conclusion.

\subsubsection{Selected half-carry windows}\label{res:selectedhalfwindow}

The selected-window construction keeps an explicit admissible word for each
terminal carry in a finite protected interval.  A supply of such windows at
all sufficiently large levels gives cofinal admissibility and hence an
infinite support with sum $1/2$:
\lword{HalfCarrySelectedWindow.lean}{831}{cofinalAdmissibility_of_selectedHalfWindows}{selected-window cofinality}
and
\lword{HalfCarrySelectedWindow.lean}{864}{exists_infinite_support_half_of_selectedHalfWindows}{the selected-window half-value conclusion}.
The unresolved input is the coherent history at each step: either a common
next-row divisor cylinder, or the canonical adjacent seam.  The theorem
states this input exactly; it does not construct the window supply and does
not prove Erd\H{o}s~\#257.

\subsubsection{A realized rewind criterion}\label{res:rewindphase}

For a realized fixed-coefficient history, the dyadic rewind phase is the
complementary binary suffix numeral:
\lword{HalfCarryRewindPhase.lean}{39}{phase_eq_denom_sub_one_sub_suffixNumeral}{the suffix-phase identity}.
Consequently the singleton criterion is exactly the inequality that this
suffix numeral covers the rest of the terminal interval,
formalised by
\lword{HalfCarryRewindPhase.lean}{65}{rewind_interval_singleton_iff_suffixNumeral_covers}{the suffix-coverage criterion}.
This converts the phase condition into support-history coordinates, but does
not prove the required coverage inequality globally.

\subsubsection{Localized protected seams}\label{res:protectedseam}

The compactness argument needs only one admissible word at cofinally many
depths.  At a localized one-hole seam, one of the two protected carries $3$
and $4$ survives; therefore a cofinal supply of realized protected seams
implies cofinal admissibility by
\lword{HalfCarryProtectedSeamConsumer.lean}{31}{cofinalAdmissibility_of_cofinalProtectedEvenSeamRealization}{the protected-seam cofinality implication},
and hence an infinite support with sum $1/2$,
\lword{HalfCarryProtectedSeamConsumer.lean}{53}{exists_infinite_support_half_of_cofinalProtectedEvenSeamRealization}{by compactness}.
The cofinal realization supply is not proved.

\subsubsection{Combining the two cofinal half-carry criteria}\label{res:cofinalwindowseam}

Assume that at each requested scale there is either a selected protected
window or a localized realized seam.  Each alternative supplies the one
admissible word needed for compactness, so this cofinal disjunction implies
an infinite support with sum $1/2$:
\lword{HalfCarryCofinalWindowOrSeamConsumer.lean}{32}{cofinalAdmissibility_of_cofinalSelectedWindowOrProtectedSeam}{by window-or-seam cofinality}
and
\lword{HalfCarryCofinalWindowOrSeamConsumer.lean}{58}{exists_infinite_support_half_of_cofinalSelectedWindowOrProtectedSeam}{by the window-or-seam half-value conclusion}.
The cofinal disjunction remains unproved.

\subsubsection{Integer-greedy first-wrap reduction}\label{res:integergreedy}

The first-wrap seam has exact truncated Mersenne weights with a quantitative
gap-dominance property.  For these weights, a small positive subset-sum defect
exists exactly when the descending integer-greedy remainder is small and
positive; this is formalised by
\lword{HalfCylinderIntegerGreedy.lean}{654}{seam_smallPositiveDefect_iff_integerGreedyRemainder}{the integer-greedy remainder equivalence}.
The remaining arithmetic input is an unbounded lower bound for that
deterministic remainder sequence.  No such lower bound is proved here.

\subsubsection{Exact finite recurrence for the greedy seam}\label{res:concreteseam}

The finite seam-word construction realizes the abstract adjacent cut by
concrete Boolean words and identifies its next word with the integer-greedy
choice:
\lword{HalfCylinderConcreteSeamAdapter.lean}{667}{seamCutNextWord_eq_greedyWord_succ}{the seam-word recurrence}.
It implements a finite recurrence, not the unbounded remainder estimate.

\subsection{Fatal gaps and the final-skip reduction}

\subsubsection{Fatal gaps and eventual right tails}\label{app:fatalright-detail}

For the concrete integer seam words, eventual right extension is equivalent
to nonmembership of $1/2$ in the Mersenne achievement set:
\lword{HalfCylinderFatalGapRightTail.lean}{787}{seamGreedyEventuallyRight_iff_half_not_mem}{the eventual-right nonmembership equivalence}.
The same condition is equivalent to the existence of a finite fatal half gap,
by
\lword{HalfCylinderFatalGapRightTail.lean}{781}{seamGreedyEventuallyRight_iff_existsFatalHalfGap}{the fatal-gap equivalence},
and also to a final skipped exponent in the real greedy orbit,
by
\lword{HalfCylinderFatalGapRightTail.lean}{802}{seamGreedyEventuallyRight_iff_exists_isLastHalfGreedySkip}{the final-skip equivalence}.
Thus this equivalence gives an exact classification of the nonmembership branch.
It neither decides whether $1/2$ belongs to the achievement set nor proves
the universal Erd\H{o}s~\#257 statement.

This is the nonmembership half of the argument summarized earlier: a last
greedy skip is equivalent to eventual right extension.  The following
paragraph records the first exceptional middle case that can be ruled out.

\paragraph{First final-middle cell.}  The first exceptional middle cell has a
separate obstruction: if the orbit extends right after the last false rank,
then the cell with transition carry $-3$ makes the lazy cofinite tail strictly
smaller than $1/2$ while its centred carry becomes negative.  This is
impossible by the analytic nonnegativity identity,
\lword{HalfCylinderFinalMiddleCellEscape.lean}{587}{finalMiddleCell_neg_three_not_last}{the middle-\(-3\) exclusion}.
It removes that one conditional terminal configuration; it does not decide
the right-tail branch or the remaining two middle cells.

The remaining global reduction is now precise.  If every genuine middle
transition other than that discharged cell has transition carry larger than
its complete future divisor-incidence tail, then an eventual right seam is
impossible and $1/2$ belongs to the Mersenne achievement set:
\lword{HalfCylinderLastProducerContradiction.lean}{387}{half_mem_mersenneAchievementSet_of_middleProducerTailEscapeExceptNegThree}{middle-tail implication}.
The stronger square-root lower-bound hypothesis is also sufficient,
\lword{HalfCylinderLastProducerContradiction.lean}{409}{half_mem_mersenneAchievementSet_of_middleProducerSqrtEscape}{square-root sufficient condition}.
Neither hypothesis is supplied here.

For finite seam prefixes the remaining tail bound sharpens further: the
future divisor-incidence tail is bounded by the cardinality of the prefix.
Thus the explicit row-scale inequality
\(|c|+\mathrm{belowPulse}+5<4\,\mathrm{remainder}\) is sufficient for
the same half-membership conclusion,
\lword{HalfCylinderMiddleCarryLowerBound.lean}{1503}{half_mem_mersenneAchievementSet_of_middleProducerCardEscape}{prefix-cardinality bound};
it is enough in turn that every late middle remainder is at least its row,
\lword{HalfCylinderMiddleCarryLowerBound.lean}{1510}{half_mem_mersenneAchievementSet_of_middleProducerRowEscape}{row-scale sufficient condition}.
These are exact conditional reductions, not proofs of the row bound.

\paragraph{Finite upper-reset band certificates.}  The companion upper
branch has a kernel-checked finite certificate: for every upper transition
from rows $13$ through $30$, every subsequent dyadic danger band is avoided.
The certificate shares each exact successor remainder across all its band
indices,
\lword{HalfCylinderUpperResetBandCertificates.lean}{78}{seamUpperResetDyadicBandEscape_through_thirty}{the verified upper-reset certificate}.
It is a bounded transition check and does not supply the cofinal band-escape
hypothesis.

\subsubsection{Positive half-membership classification}\label{app:halfmembership-detail}

The complementary branch has an equally exact seam formulation.  Membership
of $1/2$ in the Mersenne achievement set is equivalent to false successor
terminal bits beyond every bound,
\lword{HalfCylinderHalfMembershipClassification.lean}{126}{half_mem_mersenneAchievementSet_iff_unboundedTerminalFalse}{the unbounded terminal-skip equivalence};
equivalently, it is equivalent to a cofinal sequence of such false terminals,
\lword{HalfCylinderHalfMembershipClassification.lean}{203}{half_mem_mersenneAchievementSet_iff_cofinalTerminalFalse}{the cofinal terminal-skip equivalence}.
It is also equivalent to a cofinal integer-seam sequence with skipped ranks
tending to infinity,
\lword{HalfCylinderHalfMembershipClassification.lean}{213}{half_mem_mersenneAchievementSet_iff_exists_unboundedSkippedRanksAlong}{the unbounded skipped-rank equivalence}.
These equivalences isolate the required cofinal skipped-rank condition; they
do not establish that condition.

\subsubsection{Largest skipped ranks}\label{res:largestskip}

At a named largest skipped rank $d$, the lower seam word and its adjacent
upper competitor have an exact late-range gap.  In division-free form, the
identity is
\[
  3L+\bigl(3\cdot2^{s+1}+2\cdot4^{s-d}+4\bigr)=3U
\]
whenever $2s<3d$; this is formalised at
\lword{HalfCylinderLargestSkipGap.lean}{259}{three_mul_largestSkipLowerWeight_add_exactLateGap_eq_upperWeight}{the exact late-gap identity}.
The upper and middle successor branches create a terminal largest skipped
rank, while the right branch preserves an existing one:
\lword{HalfCylinderLargestSkipGap.lean}{328}{seamGreedyWord_succ_isLargestFalseRank_terminal_of_upperOrMiddle}{the terminal largest-skip branch}
and
\lword{HalfCylinderLargestSkipGap.lean}{339}{IsLargestFalseRank.seamGreedyWord_succ_of_rightBranch}{right-branch persistence}.
The first-crossing branch hypothesis required to use this identity globally
is not proved.

\subsubsection{Boundary-pulse normalization}\label{res:boundarypulse}

Above a late largest skipped rank, every filled suffix rank has zero row
pulse, so the word pulse is carried entirely by the lower prefix:
\lword{HalfCylinderBoundaryPulse.lean}{85}{exists_lowerPrefix_wordPulse_eq_of_largestFalse_late}{the lower-prefix pulse localization}.
At the first crossing of the two-thirds boundary there are exactly two cases,
$3d=2s+1$ and $3d=2s+2$, and the skipped rank has pulse respectively two or
one:
\lword{HalfCylinderBoundaryPulse.lean}{150}{rowPulse_boundary_of_late_firstCrossing}{the first-crossing pulse classification}.
This supplies the incidence normalization for a transition-carry estimate;
it does not establish that estimate.

\subsubsection{Largest-skip induction}\label{res:largestskipinduction}

Assume that at the first crossing a late largest skipped rank either
remains late at the next row or forces an upper-or-middle successor branch.
This hypothesis propagates a late largest skipped rank from row
$14$ onward,
\lword{HalfCylinderLargestSkipInduction.lean}{90}{LargestSkipLateStepSocket.largestSkipLateAt_succ}{largest-skip propagation step},
and hence yields cofinal skipped ranks and membership of $1/2$ in the Mersenne
achievement set,
\lword{HalfCylinderLargestSkipInduction.lean}{167}{half_mem_mersenneAchievementSet_of_largestSkipLateStepSocket}{largest-skip membership criterion}.
The first-crossing hypothesis itself remains unproved.

\subsubsection{The fixed-tail survival criterion}\label{res:fixedtailsocket}

A skipped rank is final exactly when its post-decision remainder lies strictly
above the complete remaining Mersenne tail:
\lword{HalfCylinderFixedTailSocket.lean}{22}{isLastHalfGreedySkip_iff_skip_and_fatal}{the final-skip criterion}.
Consequently membership of $1/2$ is equivalent to the pointwise assertion that
every actual skipped rank has remainder at most its full future tail,
\lword{HalfCylinderFixedTailSocket.lean}{73}{half_mem_iff_every_actual_skip_survives}{the half-membership survival equivalence}.
This is an exact local reformulation; the tail-survival inequality is not
proved.

\subsubsection{The carry identity at a non-right transition}\label{res:producercarry}

For a finite support aligned with a non-right transition, the residual from
$1/2$ is exactly the integer transition carry minus the complete future
divisor-incidence tail, scaled by $2^{-(2d+2)}$:
\lword{HalfCylinderProducerCarrySocket.lean}{27}{producerCarry_residual_identity}{transition-carry residual identity}.
Thus the carry exceeds that tail exactly when the support series lies below
$1/2$, and a negative carry places it above $1/2$;
\lword{HalfCylinderProducerCarrySocket.lean}{48}{producerCarry_gt_tail_iff_supportSeries_lt_half}{positive-carry criterion}
and
\lword{HalfCylinderProducerCarrySocket.lean}{145}{producerCarry_neg_forces_half_lt_supportSeries}{negative-carry criterion}.
An upper carry at most $-8$ even gives a full Mersenne-gap margin,
\lword{HalfCylinderProducerCarrySocket.lean}{190}{producerCarry_le_neg_eight_forces_gap_margin}{full-gap carry criterion}.
These are local criteria, not a proof that the required carries occur.

\subsubsection{Quarter-band endpoint cells}\label{res:quarterband}

For the adjacent seam recurrence, failure of the desired next-row upper bound
on the middle branch is exactly one multiple-of-four pulse cell,
\lword{HalfCylinderQuarterBandEndpoints.lean}{153}{seamAdjacentCut_middleNextUpperFailure_iff_pulseCell}{the middle-branch pulse-cell equivalence}.
On the terminal-reset branch, the same failure is exactly a constrained
$3\cdot2^{s-1}+k$ pulse cell,
\lword{HalfCylinderQuarterBandEndpoints.lean}{175}{seamAdjacentCut_rightNextUpperFailure_iff_pulseCell}{the reset-branch pulse-cell equivalence}.
These are arithmetic normal forms for the exceptional cells; no avoidance
statement for the concrete seam orbit is proved.

\subsubsection{Reset-deficit escape}\label{res:resetdeficit}

At a late largest skipped rank, a right successor branch pins the integer
greedy remainder between an explicit quarter threshold and carry threshold:
\lword{HalfCylinderResetDeficitEscape.lean}{200}{rightBranch_remainder_window}{the right-branch remainder window}.
Along such a branch the remainder satisfies an exact affine recurrence,
\lword{HalfCylinderResetDeficitEscape.lean}{218}{rightBranch_remainder_succ_eq}{the right-branch affine recurrence}.
Moreover, a doubled-scale deficit forces an upper-or-middle branch within the
corresponding number of rows,
\lword{HalfCylinderResetDeficitEscape.lean}{261}{upperOrMiddle_within_of_two_pow_deficit}{the deficit escape bound}.
The remaining input is the deterministic anti-concentration needed to exclude
the pinned right-branch window globally.

\subsubsection{Alignment of seam and transition carries}\label{res:seamproducer}

At a first-wrap seam, the transition carry of the left boundary support is
four times the shifted seam hole, less the explicit paired incidence pulse at
the next two rows:
\lword{HalfCylinderSeamProducerAlignment.lean}{24}{producerCarry_insert_eq_four_mul_shiftedHole_sub_pulse}{seam--carry alignment}.
This is an exact coordinate change between the seam and transition
coordinates; it does not supply a sign or size estimate for that carry.

\subsection{Finite computations, exact identities, and counterexamples}

\subsubsection{A verified initial selected window}\label{res:selectedwindowbase}

The selected half-carry construction has a checked depth-$18$ base: twelve
Boolean words cover every terminal carry in the strip $[1,12]$, with all
intermediate carries in the corresponding half strip.  The reflected table is
formalised at
\lword{HalfCarrySelectedWindowBase.lean}{79}{depth18TableCheck_eq_true}{the depth-\(18\) table verification};
the resulting selected window is defined at
\lword{HalfCarrySelectedWindowBase.lean}{129}{depth18SelectedHalfWindow}{the depth-\(18\) selected window}.
All twelve words share the same restriction at depth $13$, giving the first
next-row divisor-agreement certificate,
\lword{HalfCarrySelectedWindowBase.lean}{172}{depth18SelectedHalfWindow_nextRowDivisorAgreement}{the depth-\(18\) divisor-agreement certificate}.
This finite base does not provide selected windows at unbounded depths.

\subsubsection{Rewind-to-seam equivalence}\label{res:rewindseambridge}

For a selected window with a rewind seam, a two-value coefficient profile on
the adjacent base ancestors induces the live next-row profile,
\lword{RewindSeamOperationalBridge.lean}{140}{nextRow_profile_of_rewindSeam}{the next-row rewind profile}.
Under the stated strip and buffer inequalities, that profile realizes a
protected one-hole seam at the next row,
\lword{RewindSeamOperationalBridge.lean}{217}{protectedEvenSeamRealizedAt_succ_of_profile}{the protected-seam realization}.
The base-ancestor coefficient drop is an explicit input, not a theorem here.

\subsubsection{Selected suffix cylinders}

Put
\[
 H_N=2\lfloor\sqrt N\rfloor+4.
\]
For a Boolean word \(a=(a_0,\ldots,a_N)\), let
\[
 f_a(n)=\#\{d\leq N:d\mid n,\ a_d=1\},\qquad
 K_0(a)=1,\qquad K_{r+1}(a)=2K_r(a)-f_a(r+2).
\]

\begin{definition}[selected window and suffix cylinder]
\label{def:selected-suffix-cylinder}
A \emph{selected window} of depth \(N\) and radius \(R\) is a family
\((a^{(k)})_{1\leq k\leq R}\) of Boolean words with
\(a^{(k)}_0=a^{(k)}_1=0\) such that
\[
 1\leq K_{n-1}(a^{(k)})\leq H_n
 \quad(1\leq n\leq N),
 \qquad K_{N-1}(a^{(k)})=k.
\]
For \(M\leq N\), define the suffix numeral
\[
 \nu_{M,N}(a)=\sum_{j=M+1}^{N}a_j\,2^{N-j}.
\]
The window has a \emph{suffix cylinder at cutoff \(M\) with endpoint \(E\)}
if all \(a^{(k)}\) have the same prefix through \(M\) and
\[
 \nu_{M,N}(a^{(k)})+k=E
 \qquad(1\leq k\leq R).
\]
A \emph{profiled gap stage} additionally permits one missing carry interval
\([g_-,g_+]\subseteq[1,H_N]\): a lower suffix-cylinder family represents
\(1,\ldots,g_--1\), an upper family represents
\(g_++1,\ldots,H_N\), and the two families retain fixed consecutive
prefixes through the cutoff.
\end{definition}

\begin{proposition}[conditional suffix-cylinder feedback]
\label{res:selectsuffix}
The following statements hold.
\begin{enumerate}[label=\textup{(\roman*)},leftmargin=2.2em]
\item The verified depth-\(18\) selected window has \(R=12\) and a suffix
cylinder at cutoff \(M=13\) with endpoint \(E=17\).

\item Whenever the canonical one-row successor hypotheses hold and every
selected word has the same next coefficient \(C\), the successor remains a
suffix cylinder and its endpoint is
\[
 E'=2E-C.
\]
At a feedback row \(N+1=2(M+1)\), either the full cylinder advances or the
exceptional carry values form a profiled gap stage; the alternative includes
the protected-seam case.

\item Suppose a profiled gap stage has lower and upper next coefficients
\(C_-\) and \(C_+\), respectively.  If
\[
\begin{gathered}
 1\leq 2g_- -C_- -1,\qquad C_+\leq2g_+,\\
 2g_- -C_- -1\leq 2g_+ -C_+\leq H_{N+1},\\
 H_{N+1}+C_+\leq2H_N,
\end{gathered}
\]
then the stage advances one row with exact child gap
\[
 [g'_-,g'_+]
   =[\,2g_- -C_- -1,\;2g_+ -C_+\,].
\]
At a cutoff boundary \(N=2M+1\), \(M\geq4\), the lower and upper prefixes
acquire the forced bits \(0\) and \(1\); their consecutiveness is preserved.
Under the displayed child-gap conditions at the next row, this promoted
stage therefore advances with coefficients determined by its retained
prefixes.

\item If a two-row update swallows a protected strip of bound \(B\), and the
lower coefficient pulse is at most the upper pulse plus \(2\), then
\[
 B\leq4(g_+-g_-+1)+2. \tag{A.1}
\]
This is a necessary quarter-band condition, not a contradiction.

\item If full suffix-cylinder stages exist at arbitrarily large depths, then
there is an infinite set \(A\subseteq\N\) such that
\[
 \sum_{n\in A}\frac1{2^n-1}=\frac12.
\]
\end{enumerate}
\end{proposition}

\noindent Part~\textup{(i)} is the
\lword{SelectedSuffixCylinder.lean}{88}{depth18SelectedHalfWindow_hasSuffixCylinderAt_thirteen}{depth-\(18\) cylinder}.
Part~\textup{(ii)} is the
\lword{SelectedSuffixCylinder.lean}{136}{SelectedHalfWindow.hasSuffixCylinderAt_stepExplicit}{endpoint recurrence},
\lword{SuffixCylinderGlobalProducer.lean}{174}{CylinderStage.feedbackStep_or_protectedSeam}{feedback dichotomy},
and
\lword{SuffixCylinderInStrip.lean}{665}{CylinderStage.feedbackAdvance_or_inStripTwoSheet}{two-family gap alternative}.
Part~\textup{(iii)} is the
\lword{SuffixCylinderProfiledGap.lean}{124}{ProfiledGapStage.advanceOfNextCoeffProfile}{two-coefficient gap recurrence}
and
\lword{SuffixCylinderProfiledAutoStep.lean}{110}{ProfiledGapStage.promoteThenAdvance_hasAdjacentPrefixes}{promote-and-advance theorem}.
Part~\textup{(iv)} is the
\lword{SuffixCylinderProfiledSwallow.lean}{254}{nextBound_le_four_gapWidth_add_two_of_macro_swallow}{swallow bound}.
Part~\textup{(v)} is the
\lword{SuffixCylinderTerminalOnlyBridge.lean}{138}{exists_infinite_support_half_of_cofinalCylinderStages}{cofinal-cylinder criterion}.

\paragraph{Status and prior-art boundary.}
The proposition is a conditional propagation theorem: it does not supply the
coefficient inequalities or the cofinal family in part~\textup{(v)}.  The
strict-tail achievement-set literature supplies the surrounding Cantor-set
geometry~\cite{kovactao}, not this finite carry recurrence; the comparison
record makes no novelty claim for the recurrence itself.

\subsubsection{Verified finite suffix-cylinder stages}\label{res:prefeedback}

Before the first half-divisor feedback row, Lean constructs a full selected
window at every depth from $18$ through $27$,
\lword{SelectedSuffixCylinderPreFeedback.lean}{118}{nonempty_stage_of_eighteen_le_of_le_twentySeven}{the stages through depth \(27\)}.
At depth $27$ the common suffix cylinder remains and its endpoint covers the
whole protected strip,
\lword{SelectedSuffixCylinderPreFeedback.lean}{141}{depth27SelectedHalfWindow_hasSuffixCylinderAt_thirteen}{the depth-\(27\) suffix cylinder}
and
\lword{SelectedSuffixCylinderPreFeedback.lean}{147}{depth27SelectedHalfWindow_endpoint_covers_strip}{its strip coverage}.
The shared depth-$13$ prefix gives an explicit checked coefficient recurrence,
including the resulting depth-$27$ endpoint,
\lword{SuffixCylinderFirstFeedback.lean}{144}{stage27_endpoint_eq}{the depth-\(27\) endpoint calculation}.
This is a finite run, not an unbounded induction.
The threshold-stage continuation extends this explicit window through depth
$29$: its endpoint still covers the strip and its common suffix cylinder is
certified at every cutoff from $14$ through $25$, while cutoff $26$ is
explicitly not common,
\lword{SuffixCylinderThreshold.lean}{700}{depth29SelectedHalfWindow_endpoint_covers_strip}{the depth-\(29\) strip coverage}
and
\lword{SuffixCylinderThreshold.lean}{806}{depth29SelectedHalfWindow_not_commonRestrictionAt_twentySix}{the cutoff-\(26\) obstruction}.

An independent exact run reaches depth \(51\) with endpoint
\(51{,}327{,}745\), beyond the next feedback head threshold, and therefore
promotes the cutoff.  The promoted successor at depth \(52\) is also checked:
\lword{SuffixCylinderStage51.lean}{536}{fullCylinderStage51_feedback_headThreshold_margin}{depth-\(51\) threshold margin}
and
\lword{SuffixCylinderTerminalOnlyBridge.lean}{118}{fullCylinderStage52_halfTerminalOnlyStripWitness}{depth-\(52\) terminal-strip witness}.
These computations establish finite stages only; they do not provide the
arbitrarily large stages required by
Proposition~\ref{res:selectsuffix}\textup{(v)}.

\subsubsection{Half-divisor unit drop}\label{res:halfdivisordrop}

At the doubled feedback row, changing only the terminal half-divisor bit from
false to true increases the support coefficient by exactly one:
\lword{HalfDivisorUnitDrop.lean}{20}{supportCoeff_extend_true_eq_false_add_one_at_double}{the terminal-bit unit-drop identity}.
The same holds for any explicitly identified adjacent boundary pair,
\lword{HalfDivisorUnitDrop.lean}{35}{supportCoeff_boundaryPair_unitDrop_at_double}{the boundary-pair unit-drop identity}.
The inverse-parent recurrence transfers such a pair to the next protected
seam while retaining the exact unit-drop witness,
\lword{RewindHalfDivisorAdapter.lean}{86}{rewindBaseUnitDropAt_of_halfDivisorBoundaryPair}{the rewind unit-drop transfer}
and
\lword{RewindHalfDivisorAdapter.lean}{147}{protectedEvenSeamRealizedAt_succ_of_fullStrip_rewindHalfDivisorBoundaryPair}{the protected-seam transfer}.
This is the local coefficient input for the rewind-seam profile, not a global
construction of such boundary pairs.

\subsubsection{A finite rewind-boundary obstruction}\label{res:rewindboundarycounterexample}

The scalar one-row rewind data alone still does not supply that missing
boundary predicate.  Lean gives a concrete depth-$26$ selected two-word
window whose depth-$27$ explicit step has rewind history $[1]$, is a width-two
seam pair, and satisfies the doubled-row identity, but has no depth-$13$
half-divisor boundary-pair witness:
\lword{RewindBoundaryPairCounterexample.lean}{169}{depth27_actualRewindSeam_not_boundaryPair}{the depth-\(27\) boundary-pair counterexample}.
Thus the protected-window buffer in the compactness criterion is a genuine
hypothesis, not a consequence of the scalar seam facts alone.

\subsubsection{Explicit carry coordinates at the final non-right transition}
\label{res:producercarrycoordinates}

For every adjacent cut at rank $s\ge5$, the transition carry of the
terminal-augmented below support equals the concrete middle coordinate:
\[
  4\,\mathrm{remainder}-\mathrm{belowPulse}-4,
\]
and identifies the terminal-augmented above support with the negative upper
coordinate
\[
  -\bigl(4\,\mathrm{overshoot}+\mathrm{abovePulse}+4\bigr).
\]
The generic row-word identity is
\lword{HalfCylinderProducerLowerBound.lean}{105}{producerCarry_insert_seamWordSupport_eq}{row-word transition identity},
with the below and above specializations at
\lword{HalfCylinderProducerLowerBound.lean}{150}{producerCarry_insert_seamBelowSupport_eq_middleCoordinate}{middle-coordinate identity}
and
\lword{HalfCylinderProducerLowerBound.lean}{183}{producerCarry_insert_seamAboveSupport_eq_neg_upperCoordinate}{upper-coordinate identity}.
These identities make the remaining quantitative task exact: they do not
prove exponential lower bounds, sign persistence, or a universal escape
mechanism.

\subsubsection{Strategy boundary: Campbell shift synchronization}
\label{res:campbellboundary}

A proposed route combines a finite Mersenne shift from a putative last greedy
skip with Campbell's quarter-exponent prime progression.  The exact parameter
requirements cannot be synchronized: if the combined progression modulus is
$d\ge2$, the phase factor is at most $d$, the fatal-window prime cap is at
most four times that phase, and Campbell's applicability requires
$d^4$ below the cap, then the inequalities are inconsistent.  Formalised:
\lword{CampbellShiftSynchronization.lean}{299}{campbellQuarterExponent_periodFreeze_noSync}{the period-freeze incompatibility}
and
\lword{CampbellShiftSynchronization.lean}{320}{campbellPhaseFreeze_parameters_impossible}{the parameter incompatibility}.

The associated shifted-zero condition is not an independent hypothesis.  It
is exactly equivalent both to infinitely many actual greedy skips and to
$1/2\in\mathcal A$:
\lword{CampbellShiftSynchronization.lean}{471}{shiftWindowZero_iff_greedyMersenneSkippedSupport_infinite}{the shifted-zero/skip equivalence}
and
\lword{CampbellShiftSynchronization.lean}{488}{shiftWindowZero_iff_half_mem_mersenneAchievementSet}{the shifted-zero/half-membership equivalence}.
This closes one synchronization strategy and one reformulation; it does not
advance either endpoint beyond the already stated half-achievement problem.

}

\subsection{What remains open}

The two directions remain separate.
\begin{enumerate}[label=\textup{(\alph*)},leftmargin=2.2em]
\item The universal problem is to prove \(X_A\notin\Q\) for every infinite
\(A\subseteq\N\setminus\{0\}\).  Theorem~\ref{res:support} proves several
families, and Theorem~\ref{res:rigidity-main} constrains a hypothetical
rational support, but neither result supplies the universal quantifier.
\item The counterexample problem is to prove \(1/2\in\mathcal A\).
Theorem~\ref{res:halfmembership} identifies this with infinitely many greedy
skips, while Theorem~\ref{res:fatalright} reduces the current final-skip
argument to two exact obligations: exclude the remaining exceptional cells
\(C_D=-2,-1\), and prove the tail inequality~\textup{(4.2)} at every other
non-\(-3\) middle transition.  Neither obligation is discharged.
\end{enumerate}
Appendix~\ref{app:half-technical} records the exact local reductions, finite
certificates, counterexamples to stronger local assertions, and closed
strategy branches supporting this conclusion.  They refine the obstruction;
they do not constitute a second theorem spine.

\section{The totient constant $S$ (Erd\H{o}s \#249)}\label{sec:249}

Whether
\[
S=\sum_{n\ge1}\frac{\ph(n)}{2^n}
\]
is irrational remains open.  The argument has four logically distinct
layers.  First, a finite Farey computation excludes every rational
denominator up to an explicit bound.  Second, irrationality is characterised
exactly by the existence of finite certificates against every possible
eventual binary period.  Third, many such certificates are verified at
bounded parameters.  Fourth, and still missing, is a theorem supplying them
at unbounded parameters.  Only the fourth layer would settle the problem.

The main section states this theorem spine.  Appendix~\ref{app:totient-technical}
contains the cone, projection, denominator, parity, and closed-strategy
reductions that refine the open fourth layer.

\subsection{Unconditional: no small denominator}\label{sec:uncond}

\begin{theorem}[denominator exclusion]\label{res:farey}
Set
\[
\Qzero:=79\,639\,646\,646\,701\,375\,323\,355\,774\,875\,831\,053
       \approx7.96\times10^{34}.
\]
For every integer $p$ and every $q$ with $1\le q\le\Qzero$,
\[
S \;\ne\; \frac{p}{q}.
\]
Equivalently, if $S$ is rational then its reduced denominator exceeds that bound.
\end{theorem}
\noindent Formalised: \lloc{CertificateKernel.lean}{18055}.
The bound is reached by a ladder of finite, elementary exclusions,
\[
4838 \;\to\; 4\,194\,304 = 2^{22} \;\to\; 2.49\times10^{17}\ (\text{Farey } K{=}120)
\;\to\; 7.96\times10^{34}\ (\text{Farey } K{=}240),
\]
each a mediant/Farey argument that traps any candidate $m/q$ inside a forbidden
gap. The mediant framework is historically associated with Farey's 1816 note
on vulgar fractions~\cite{farey1816}; the finite exclusions themselves are
Lean-checked arguments in the formal-source checkpoint. See
\lword{CertificateKernel.lean}{15347}{tsum_totient_div_pow_two_ne_ratCast_of_den_le_4838}{denominator exclusion through 4838}
for the first rung and \lword{GapFareyBound.lean}{88}{gap_check_window_1_120_le_248672326362367909}{120-term Farey bound},
\lword{GapFareyBound.lean}{176}{gap_check_window_1_240_le_79639646646701375323355774875831053}{240-term Farey bound}
for the two Farey windows.

This is an explicit lower bound on the denominator of any rational
representation, not a proof of irrationality. A rational with a larger
denominator is not excluded by it.

\begin{remark}[one bound, four representations]
Because the ladder identities of Section~\ref{sec:ladder} are equalities, the same
finite record transfers verbatim to three other representations of the constant:
to the positive lift $\sum_d A(d)/(2^d-1)$
(\lloc{CertificateKernel.lean}{18144}),
to the signed M\"obius-square constant $T=\sum_d\mu(d)/(2^d-1)^2$ at half the bound
(\lloc{CertificateKernel.lean}{18158}),
and to the coprimality probability of Section~\ref{sec:geom} at half the bound.
Half appears because $S=\tfrac12+T$ turns a denominator $d$ for $T$ into $2d$ for $S$.
\end{remark}

\subsection{The geometric picture: $S$ as coprime-pair mass}\label{sec:geom}

Let $X,Y$ be independent fair-coin waiting times, $\Pr(X=n)=2^{-n}$ for $n\ge1$.
Then $\Pr(d\mid X)=1/(2^d-1)$, so the Lambert transform is a divisor calculus of
$X$: $L(f)=\mathbb E[(f*\zeta)(X)]$. The rungs read as $L(\mu)=\Pr(X=1)=1/2$,
$L(\ph)=\mathbb E[X]=2$, $L(\mathbf 1)=\mathbb E[\tau(X)]=E$, and $L(A)=\mathbb E[\ph(X)]=S$.

Counting visible lattice points $(a,b)$ with $a+b=n$, $a>0$, $\gcd(a,b)=1$ gives
exactly $\ph(n)$, uniformly in $n$, so $S$ is itself a coprime-pair mass.

\begin{proposition}[coprimality probability]\label{res:coprime}
\[
S \;=\; \tfrac12 \;+\; \Pr\!\big(\gcd(X,Y)=1\big)
\;=\; \tfrac12 \;+\!\!\sum_{\substack{a,b\ge1\\ \gcd(a,b)=1}}\!\! 2^{-(a+b)}.
\]
\end{proposition}
\noindent Formalised: \lword{CertificateKernel.lean}{18229}{totient_series_eq_half_add_visible_coprime_pairs}{coprime-pair representation},
with the underlying lattice identity at
\lword{CertificateKernel.lean}{18216}{tsum_visible_coprime_pairs_eq_totient_series}{visible-lattice-point identity}.
Base $2$ is the one point where $\Pr(X=a)\Pr(Y=b)=2^{-(a+b)}$ needs no normalising
constant, so the totient sum and the coprimality probability line up with no
boundary correction. In this reading, Erd\H{o}s \#249 asks whether two independent
fair-coin waiting times are coprime with irrational probability.
After reducing $(X,Y)$ by their gcd, the same distribution has an exact
Stern--Brocot cylinder decomposition.  Appendix~\ref{app:lambert-complements}
states the resulting probability law and the sharp Fibonacci stability of its
alternating runs.  These are unconditional structural theorems about the
representation of $S$; they do not supply the denominator survival needed for
irrationality.

\subsection{Exact certificate characterisation}\label{sec:reduction}

The reduction uses only the elementary bound $\ph(n)\le n$, geometric tail
estimates, and finite integer arithmetic.
For $t\ge1$, write $M_t=\lcm(1,\ldots,t)$.

\begin{definition}[tail differences and finite certificates]
\label{def:totient-certificate}
For $N\ge0$, define the scaled tail
\[
R_N \;=\; \sum_{m\ge1}\frac{\ph(N+m)}{2^m}.
\]
For $h\ge1$ and $L\ge1$, define
\[
\Delta_{h,N,L}=\sum_{j=0}^{L-1}\big(\ph(N+h+1+j)-\ph(N+1+j)\big)\,2^{\,L-1-j}\in\Z,
\]
and say that $(h,N,L)$ is a \emph{finite certificate} when
\[
\textsf{certifiedKill}(h,N,L):\quad
N+h+L+2 \;<\; \big(\Delta_{h,N,L}\bmod 2^L\big) \;<\; 2^L-(N+h+L+2).
\]
\end{definition}

The integer $\Delta_{h,N,L}$ contains the first $L$ binary places of
$R_{N+h}-R_N$; the two margins bound the omitted tail using
$\ph(n)\le n$.  Thus the predicate is decidable by finite integer arithmetic,
while its conclusion concerns the infinite tail difference.

\begin{theorem}[exact tail-period characterisation]\label{res:red21}
$S$ is irrational if and only if, for every period $h\ge1$ and every threshold
$N_0$, there exist $N\ge N_0$ and $L$ with $\textsf{certifiedKill}(h,N,L)$.
\end{theorem}
\noindent Formalised: \lword{LcmConeFlatness.lean}{412}{irrational_totient_series_iff_certificate_supply}{certificate-supply equivalence};
the tail $R_N$ and the predicate are
\lword{TotientTailPeriodKiller.lean}{57}{totientTail}{scaled totient tail} and
\lword{TotientTailPeriodKiller.lean}{72}{certifiedKill}{finite certificate predicate}.

Every multiple of an eventual period is again a period.  The least common
multiple $M_t$ therefore provides a useful sufficient diagonal form of
Theorem~\ref{res:red21}.

\begin{corollary}[diagonal reduction]\label{res:diag}
Suppose that for every threshold $t_0$ there are parameters $t\ge t_0$ and
$L$ for which $\textsf{certifiedKill}(M_t,M_t,L)$ holds. Then $S$ is
irrational.
\end{corollary}
\noindent Formalised: \lword{LcmDiagonalReduction.lean}{112}{irrational_totient_series_of_lcm_diagonal_certificate_supply}{diagonal-supply criterion}.

\begin{proposition}[certificate completeness]\label{res:complete}
For every $h\ge1$ and $N\ge0$,
\[
\bigl(\exists L\ge1,\ \textsf{certifiedKill}(h,N,L)\bigr)
\quad\Longleftrightarrow\quad
R_{N+h}-R_N\notin\Z.
\]
Moreover, the following are equivalent:
\[
\begin{aligned}
 &S\notin\Q,\\
 &R_{N+h}-R_N\notin\Z
      &&\text{for every }h\ge1\text{ and }N\ge0,\\
 &\exists L\ge1,\ \textsf{certifiedKill}(h,N,L)
      &&\text{for every }h\ge1\text{ and }N\ge0.
\end{aligned}
\]
\end{proposition}
\noindent Formalised:
\lword{LcmConeFlatness.lean}{316}{exists_certifiedKill_iff_tail_diff_notMem_int}{certificate-completeness equivalence},
\lword{LcmConeFlatness.lean}{386}{irrational_totient_series_iff_all_tail_diffs_nonintegral}{all-differences criterion},
and
\lword{LcmConeFlatness.lean}{399}{irrational_totient_series_iff_pointwise_certificates}{pointwise certificate criterion}.
The converse direction uses the fact that integrality of even one
positive-shift tail difference forces $S$ to be rational
(\lword{LcmConeFlatness.lean}{357}{not_irrational_of_tail_diff_mem_int}{integral-difference rationality criterion}).
Thus the finite predicate is complete at each fixed pair $(h,N)$; the open
problem is to prove that the required pairs occur without bound.

\begin{example}[verified finite range]\label{res:deposits}
Lean checks:
\begin{itemize}[leftmargin=1.4em]
\item $\textsf{certifiedKill}(h,12,16)$ for every $1\le h\le8$;
\item a fixed-window certificate for every $1\le h\le16$;
\item diagonal certificates $\textsf{certifiedKill}(M_t,M_t,L)$ at
\[
t\in\{1,2,3,4,5,7,8,9,11,13,16,17,19,23,25,27,29,31,32,37,41,43,
47,49,53,59,61,64\};
\]
\item a finite joint certificate refuting cone flatness at a tabulated family
of vertices.
\end{itemize}
\end{example}
\noindent Formalised:
\lword{TotientTailPeriodKiller.lean}{404}{certifiedKill_all_small}{small-window certificate family},
\lword{CarrySurvivorExtinction.lean}{574}{certifiedKill_all_upto_sixteen}{periods through sixteen},
\lword{LcmDiagonalReduction.lean}{237}{certifiedKill_periodLcm_diagonal_upto_six}{diagonal certificates through scale six},
\lword{DiagonalPincerCertificatesT64.lean}{1967}{certifiedKill_diagonal_all_imported_through_t64}{imported diagonal certificates through scale 64},
and
\lword{LcmConeNonflat.lean}{506}{coneNonflatCert_cells}{cone-cell certificate table}.

For example, at $(h,N,L)=(1,12,16)$,
\[
\Delta_{1,12,16}=-143140,\qquad
\Delta_{1,12,16}\bmod 2^{16}=53468,
\]
while the omitted-tail allowance is $12+1+16+2=31$.  Hence
\[
31<53468<2^{16}-31,
\]
which is exactly $\textsf{certifiedKill}(1,12,16)$ and therefore proves
$R_{13}-R_{12}\notin\Z$.  Every row above has this exact status: it proves a
particular non-integral tail difference, not an unbounded supply.

\subsection{What remains open}\label{sec:249-open}

By Theorem~\ref{res:red21}, the unresolved statement is precisely
\[
\forall h\ge1\ \forall N_0\ge0\ \exists N\ge N_0\ \exists L\ge1,\qquad
\textsf{certifiedKill}(h,N,L).
\tag{5.1}\label{eq:unbounded-certificate-supply}
\]
The diagonal hypothesis in Corollary~\ref{res:diag} is a sufficient
one-parameter strengthening of this statement.  The verified scales in
Example~\ref{res:deposits} form a finite set and do not establish either
quantifier pattern.  No theorem in the formal source proves
\eqref{eq:unbounded-certificate-supply}; consequently no irrationality claim
for $S$ is made here.

Appendix~\ref{app:totient-technical} records exact reformulations of this open
obligation, finite projections and denominator formulae, obstructions to
several stronger local criteria, and closed strategy branches.  They sharpen
the target but do not supply its missing unbounded quantifier.

\long\def\TotientTechnicalAppendix{%

\subsection{Exact normal forms for the open supply}

\paragraph{Cone flatness and completeness.} On the cone
$\{kM_t:k\ge1\}$, rationality forces a single fractional constant across the
whole cone.

\begin{proposition}[cone flatness]\label{res:flat}
If $S$ is rational, then for every sufficiently large $t$ there is
$\theta_t\in[0,1)$ such that the fractional part of the tail $R_{kM_t}$ equals
$\theta_t$ for every $k\ge1$.
\end{proposition}
\noindent Formalised: \lword{LcmConeFlatness.lean}{90}{rational_totient_series_forces_lcm_cone_flatness}{cone flatness under rationality}.

The certificates are exact: a certificate exists exactly
when the object it certifies is non-integral. Thus the finite computations are
machine-checked witnesses rather than numerical experiments.

\begin{proposition}[certificate completeness, expanded form]\label{app:complete-detail}
For all $h,N$,
\[
\big(\exists L,\ \textsf{certifiedKill}(h,N,L)\big)
\quad\Longleftrightarrow\quad
R_{N+h}-R_N\notin\Z.
\]
\end{proposition}
\noindent Formalised: \lword{LcmConeFlatness.lean}{316}{exists_certifiedKill_iff_tail_diff_notMem_int}{certificate-completeness equivalence}.
Conversely, integrality of one positive-shift tail difference forces $S$ to be
rational.  Thus $S$ is irrational if and only if every positive-shift tail
difference is non-integral, equivalently if every such pair admits a finite
certificate.  These normal forms are formalised at
\lword{LcmConeFlatness.lean}{357}{not_irrational_of_tail_diff_mem_int}{integral-difference rationality criterion},
\lword{LcmConeFlatness.lean}{386}{irrational_totient_series_iff_all_tail_diffs_nonintegral}{all-differences criterion},
and \lword{LcmConeFlatness.lean}{399}{irrational_totient_series_iff_pointwise_certificates}{pointwise certificate criterion}.
Finite certificates witness non-integrality. A joint refuter
can interrogate a menu of cone vertices at once and is sharper than any pairwise
test: it refutes the single-fractional-part model that rationality predicts on a
finite sample of cone vertices
(\lword{LcmConeNonflat.lean}{126}{exists_nonintegral_pair_of_coneNonflatCert}{finite cone-nonflatness refuter}).

\paragraph{A reduced-denominator finite predicate.}  Separately, there is a
stricter finite condition for an odd reduced denominator
$u$.  Its predicate $\textsf{ReducedDenominatorUnitGapCert}(u,N,K)$ asks that
every candidate integer in the associated dyadic tail interval be non-coprime
to $u$.  At a prime power $u=p^e$, the condition is equivalent to the exact
alternative that the interval has no candidate, or that it has one candidate
and $p$ divides that candidate.  Thus such an interval has at most one
candidate.

\begin{proposition}[prime-power unit-gap ceiling]\label{res:unitgap}
For every prime power $p^e$ with $e>0$, a
$\textsf{ReducedDenominatorUnitGapCert}(p^e,N,K)$ has candidate count at most
one.  The row $(u,N,K)=(3,3,5)$ satisfies this predicate, although its
corresponding ordinary empty-gap inequality fails.
\end{proposition}
\noindent Formalised:
\lword{PrimitiveWeightCertificate.lean}{54}{reducedDenominatorUnitGapCert_prime_pow_iff}{prime-power unit-gap characterisation},
\lword{PrimitiveWeightCertificate.lean}{111}{reducedDenominatorCandidateCount_prime_pow_le_one}{prime-power one-candidate bound},
\lword{PrimitiveWeightCertificate.lean}{151}{reducedDenominatorUnitGapCert_three_window_3_5}{three-window unit-gap certificate}, and
\lword{PrimitiveWeightCertificate.lean}{157}{not_ordinary_gap_certificate_three_window_3_5}{three-window ordinary-gap failure}.
This is a finite certificate-level strengthening.  We do not establish that
rationality of $S$ supplies these predicates at unbounded parameters.

\paragraph{The certificates are non-empty.} The initial segment of the required
supply is machine-checked, so the reduction is not vacuous. By completeness, each
of these is a kernel-checked proof of a specific non-integral tail difference,
not a numerical experiment.

\begin{example}[verified finite instances, expanded record]\label{app:deposits-detail}
The following are checked by the Lean kernel:
\begin{itemize}[leftmargin=1.4em]
\item $\textsf{certifiedKill}(h,12,16)$ for every period $h\in[1,8]$;
\item certified kills for every period $h\in[1,16]$ at a fixed window;
\item diagonal certificates $\textsf{certifiedKill}(M_t,M_t,L)$ at the 28 scales
  $t\in\{1,2,3,4,5,7,8,9,11,13,16,17,19,23,25,27,29,31,32,37,41,43,
  47,49,53,59,61,64\}$;
\item a tabulated joint cone-vertex certificate.
\end{itemize}
\end{example}
\noindent Formalised: \lword{TotientTailPeriodKiller.lean}{404}{certifiedKill_all_small}{small-window certificate family},
\lword{CarrySurvivorExtinction.lean}{574}{certifiedKill_all_upto_sixteen}{periods through sixteen},
\lword{LcmDiagonalReduction.lean}{237}{certifiedKill_periodLcm_diagonal_upto_six}{diagonal certificates through scale six}
(with \lword{LcmDiagonalReduction.lean}{242}{certifiedKill_periodLcm_diagonal_at_seven}{scale-seven diagonal certificate},
\lword{LcmDiagonalReduction.lean}{246}{certifiedKill_periodLcm_diagonal_at_eight}{scale-eight diagonal certificate}),
\lword{DiagonalPincerCertificatesT64.lean}{1967}{certifiedKill_diagonal_all_imported_through_t64}{imported diagonal certificates through scale 64},
and \lword{LcmConeNonflat.lean}{506}{coneNonflatCert_cells}{cone-cell certificate table}. These instances show
that each layer of the reduction is non-vacuous and, by completeness, exact. They
are not evidence that the unbounded supply exists.

The first row of this finite table can be read without any search procedure.
Take $(h,N,L)=(1,12,16)$.  The integer in the definition is
\[
 \Delta_{1,12,16}
 =\sum_{j=0}^{15}\bigl(\ph(14+j)-\ph(13+j)\bigr)2^{15-j}
 =-143140.
\]
Its least nonnegative residue modulo $2^{16}$ is $53468$.  The tail allowance
is $N+h+L+2=31$, and hence
\[
       31<53468<65536-31=65505.
\]
This is precisely $\textsf{certifiedKill}(1,12,16)$.  The soundness theorem
(\lword{TotientTailPeriodKiller.lean}{262}{tail_diff_notMem_int_of_certifiedKill}{certificate soundness})
therefore gives $R_{13}-R_{12}\notin\Z$.  It is one exact non-integral tail
difference, not an assertion that such differences occur at unbounded scales.

\paragraph{An exact target interval.} The diagonal calculation also admits a
more local normal form.  Put $H_t=M_t$ and
\[
 D_t=R_{2H_t}-R_{H_t}.
\]
Separate an exact rational M\"obius shadow
$Q_t=a_t/d_t$ in lowest terms and define the remaining \emph{foreign defect}
by $F_t=D_t-Q_t$.  The full target is the residue class of $-a_t$ modulo $d_t$:
\[
 \textsf{FullTargetHit}(t)
 \quad\Longleftrightarrow\quad
 \exists k\in\Z,\qquad d_tF_t=-a_t+d_tk.
\]
Since $D_t=a_t/d_t+F_t$, this is exactly the condition $D_t\in\Z$; no
complementary term is discarded or assigned a conjectural sign.

\begin{proposition}[diagonal pincer decomposition]\label{res:pincer}
For every $t$, diagonal integrality is equivalent to the foreign diagonal
defect meeting the full target.  Consequently, full-target avoidance at
unbounded scales implies irrationality of $S$.
\end{proposition}
\noindent Formalised:
\lword{DiagonalPincerDecomposition.lean}{215}{diagonal_int_iff_foreignDiagonalDefect_hits_fullTarget}{full-target equivalence}
and
\lword{DiagonalPincerDecomposition.lean}{290}{irrational_totientSeries_of_full_target_avoidance_supply}{unbounded target-avoidance criterion}.
The proposition is an exact reduction; the unbounded avoidance hypothesis is
still open.

Two algebraic inputs make this interval test explicit.  First, truncate the
M\"obius-square identity of Proposition~\ref{res:lens} after $D$ terms.  The
resulting rational quantity gives a Lambert-projected centre $C_{H,D}$ for the
scaled diagonal difference.  The error is exactly a signed squared-Mersenne
tail:
\[
 \textsf{scaledDiagonal}(H)-C_{H,D}
   =c_H\sum_{d>D}\frac{\mu(d)}{(2^d-1)^2},
\]
where $c_H$ is the explicit diagonal coefficient in the formal source.  The
absolute value of the sum is bounded geometrically by
$4/[3(2^{D+1}-1)^2]$.  Thus a rational separation of the centre from every
integer by more than the scaled bound proves target avoidance; the argument
does not estimate the sign of the tail.

\begin{proposition}[squared-Mersenne enclosure]\label{res:squareenclosure}
The scaled diagonal tail difference differs from its Lambert-projected centre
by the squared-Mersenne tail.  Separation of the centre from the full target by
more than the proved tail bound forces target avoidance.
\end{proposition}
\noindent Formalised:
\lword{SquaredMersenneDiagonalEnclosure.lean}{116}{scaleDiagonalTailDifference_sub_lambertProjectedDiagonal}{squared-Mersenne tail identity}
and
\lword{SquaredMersenneDiagonalEnclosure.lean}{237}{scaleFullTarget_miss_of_lambert_projected_separation}{Lambert-projected separation criterion}.

\subsection{Finite projections and verified instances}

\paragraph{Finite foreign-residue projection.}  There is a second finite
presentation of the same obstruction.  For a cutoff $D$, split the finite
M\"obius residue diagonal into the channels $d\leq D$ which divide $H$ and
the remaining channels.  The divisor part is exactly the truncated M\"obius
shadow, while every retained nondivisor channel is explicit.  Above the
stable cutoff $2H$, a finite window of omitted channels is bounded by
\[
 c_H\left(\frac{2}{2^D}+\frac{4}{3\,4^D}\right),
\]
where $c_H$ is the diagonal coefficient.  Thus, if the actual foreign defect
agrees with this projection within the displayed bound, and the projected
quantity is farther than that bound from every integer, then the full target
is missed.

\begin{proposition}[finite foreign-residue projection]\label{res:foreignresidue}
The finite residue diagonal splits exactly into projected foreign and divisor
channels.  A controlled projected foreign defect with the stated integer
separation cannot meet the full target.
\end{proposition}
\noindent Formalised:
\lword{ActualForeignResidueProjection.lean}{307}{finiteResidueDiagonal_eq_projectedForeign_add_divisor}{foreign--divisor channel split},
\lword{ActualForeignResidueProjection.lean}{275}{abs_foreignTailWindow_le_foreignComplementBound}{omitted-channel bound},
and
\lword{ActualForeignResidueProjection.lean}{348}{scaleFullTarget_miss_of_projected_separation}{projected-separation criterion}.
The passage from finite windows to the analytic foreign complement, including
the corresponding kernel identity, is not supplied here.  Nor is an
unbounded family of separated projections; these are parts of the open
certificate-supply problem.

Second, the rational M\"obius shadow has an exact reduced denominator.  If
$r_t$ is the square-free kernel of $H_t$, $s_t=H_t/r_t$, and $J(r_t)$ is the
odd Jordan scalar defined in the algebraic shadow chain, then
\[
 d_t=
 \frac{2^{r_t}-1}
 {\gcd\!\bigl(2^{r_t}-1,\,s_t\lvert J(r_t)\rvert\bigr)}.
\]
For $t\geq5$, the product of the Mersenne factors $2^p-1$ over primes
$t/2<p\leq t$ divides $d_t$; in particular $d_t\geq2^{t/2}$.  This is a
denominator theorem for $Q_t$ only.  Cancellation by $F_t$ is exactly what the
full-target condition retains.

\begin{proposition}[Mersenne-shadow denominator]\label{res:shadowden}
At lcm height $t$, the reduced denominator of the scaled M\"obius shadow is
given exactly by the formal denominator formula.  The upper-half Mersenne
product divides it and yields the stated lower bounds.
\end{proposition}
\noindent Formalised:
\lword{MersenneShadowDenominatorGrowth.lean}{99}{exists_upperHalf_mersenne_channel_dvd_den_with_bounds}{upper-half Mersenne divisor bound}
and
\lword{MersenneShadowDenominatorGrowth.lean}{147}{lcmHeight_scaledMobiusShadow_den_exact}{exact shadow denominator}.

Finally, the remaining diagonal error can be represented by a finite integer
projection and then by binary suffix data.  At depth $L$, define
\[
 W_{t,L}=P_L(2H_t)-P_L(H_t),\qquad
 \rho_{t,L}=W_{t,L}\bmod 2^L,
\]
where $P_L(M)$ is the length-$L$ window numerator underlying
Definition~\ref{def:totient-certificate}.  The finite
residue certificate is the asymmetric safe interval
\[
 H_t+L+2<\rho_{t,L}<2^L-(2H_t+L+2).
\]
The two unequal margins are the tail bounds at the two cone vertices $H_t$ and
$2H_t$; replacing them by a symmetric slogan would lose part of the checked
statement.

\begin{corollary}[finite residue projection]\label{res:freshloss}
The diagonal residue certificate is equivalent to its projection
inequality, and the projection forces full-target avoidance.  Hence an
unbounded projection supply, or either stronger adjacent-suffix supply stated
in the corresponding theorem, implies irrationality of $S$.
\end{corollary}
\noindent Formalised:
\lword{DiagonalFreshLossBridge.lean}{61}{diagonalFreshLossResidueCert_iff_projectionCert}{residue--projection equivalence}
and
\lword{DiagonalFreshLossBridge.lean}{191}{irrational_totientSeries_of_diagonalFreshLossProjectionSupply}{unbounded projection criterion}.
Finite instances of this implication are proved, but no unbounded projection
or suffix supply is established.

\paragraph{Primality certificates through $t=64$.}\label{res:t64closure}
A finite family of Lucas certificates proves the primality facts required by
the diagonal-pincer certificates through $t=64$; for example,
\lword{DiagonalPincerPrimeCertificates/ClosureT64.lean}{51759}{prime_lucas_2364533768205644535022723273}{representative Lucas primality certificate}.
This is a finite theorem family, not an unbounded certificate supply or an
irrationality proof.

\paragraph{The diagonal certificate at $t=64$.}\label{res:t64endpoint}
Using those primality facts, Lean checks the individual diagonal-pincer
certificate
at $t=64$,
\lword{DiagonalPincerCertificateT64Endpoint.lean}{1928}{certifiedKill_diagonal_t64}{scale-64 diagonal certificate}.
This is one finite endpoint certificate, not an unbounded supply.

\paragraph{A three-rank odd window.}\label{res:flexibleoddwindow}
At an analytically admissible odd power-of-two rank, centrality at any one of
the canonical odd rank or its next two base-four successors yields the
required adjacent-suffix gap condition,
\lword{DiagonalFlexibleOddWindowSupply.lean}{93}{diagonalAdjacentSuffixGapSupply_of_powerTwoOddGuardThreeRankBand}{three-rank suffix-gap implication}.
Thus a cofinal three-rank band supply implies irrationality,
\lword{DiagonalFlexibleOddWindowSupply.lean}{136}{irrational_totientSeries_of_powerTwoOddGuardThreeRankBandSupply}{cofinal three-rank criterion};
that supply is not proved.

\paragraph{Affine form of the three-rank condition.}\label{res:flexibleoddaffine}
With the stated short-correction bounds, the three central-band predicates are
equivalent to one three-scale affine escape condition,
\lword{DiagonalFlexibleOddWindowAffine.lean}{236}{threeRankBandAt_iff_threeScaleAffineEscapeAt_of_shortCorrections}{three-scale affine equivalence}.
A zero-start affine orbit shows that a uniform correction envelope alone does
not force escape at any finite number of ranks,
\lword{DiagonalFlexibleOddWindowAffine.lean}{225}{zeroAffineOrbit_stays_in_scaled_edge}{zero-orbit obstruction}.

\subsection{Fresh-prime and bit-level boundaries}

\paragraph{Fresh-prime deficit decomposition.}  The literal foreign channel
also has an endpoint form suited to curvature tests.  Retain only the prime
support already present at $H$ in the two endpoints $H+s$ and $2H+s$, and
write the resulting discrepancy from the actual totient as an endpoint
deficit.  Each such deficit is nonnegative.  The actual diagonal increment
then equals the old-prime increment plus the lower deficit minus the upper
deficit.  The second and five-point curvature operators preserve this exact
split; the latter has coefficient row $[-8,4,2,1,1]$.  Consequently, the old
curvature margin loses at most the stated adverse weighted deficit.

\begin{proposition}[fresh-prime deficit decomposition]\label{res:freshprimedeficit}
The endpoint fresh deficits are nonnegative, give an exact decomposition of
the foreign channel, and bound the loss from the old-prime curvature margins.
\end{proposition}
\noindent Formalised:
\lword{FreshPrimeDeficitDecomposition.lean}{91}{endpointFreshDeficit_nonneg}{endpoint-deficit nonnegativity},
\lword{FreshPrimeDeficitDecomposition.lean}{184}{finiteForeignChannelIncrement_eq_endpointFreshDeficit_sub}{foreign-increment decomposition},
and
\lword{FreshPrimeDeficitDecomposition.lean}{269}{oldFivePoint_sub_secondBranchAdverseDeficit_le_actual}{adverse-deficit curvature bound}.
This supplies the exact correction term for the finite evaluations.  It supplies
neither a cofinal deficit bound nor an irrationality criterion.

\paragraph{One-bit lifting.}  The power-of-two signed-margin condition admits a
separate exact arithmetic reduction.  If a signed difference is already
divisible by $2^k$, lifting its congruence to $2^{k+1}$ is equivalent to
evenness of the quotient after division by $2^k$.  Failure of the lift gives
the unique translated residue class.  In particular, a displayed difference
$P-N=16c$ agrees modulo $32$ exactly when $c$ is even; otherwise the residue
is the offset class by $16$.

\begin{proposition}[one-bit signed-margin lift]\label{res:powertwobitlift}
The fifth-bit question after fourth-bit agreement is exactly one cofactor
parity test, with the complementary offset class characterising failure.
\end{proposition}
\noindent Formalised:
\lword{PowerTwoBitLift.lean}{34}{two_pow_succ_dvd_iff_quotient_even}{one-bit divisibility lift},
\lword{PowerTwoBitLift.lean}{111}{modEq_thirtyTwo_iff_even_difference_cofactor}{mod-32 cofactor criterion},
and
\lword{PowerTwoBitLift.lean}{129}{not_modEq_thirtyTwo_iff_offset_sixteen}{offset-16 complement}.
The centered-state formulation identifies the signed-margin interval
condition with failure of this next lift, subject to its stated correction and margin
budget:
\lword{PowerTwoCenteredBitLift.lean}{121}{actualPenultimateSignedMargin_centeredLift_iff_not_twice_dvd}{centred-lift equivalence}
and
\lword{PowerTwoCenteredBitLift.lean}{159}{actualPenultimateSignedMargin_centeredThirtyTwo_iff_not_thirtyTwo_dvd}{centred mod-32 equivalence}.
This is a finite arithmetic equivalence.  It does not supply the required
cofactors at unbounded scales and gives no irrationality conclusion.

\paragraph{Parity is not a margin certificate.}  The preceding lift does not
settle the reduced signed-margin interval condition.  For every positive even
modulus and every radius in the stated nontrivial range, each parity class
contains both a residue inside the interval and one outside it.  The same
statement applies to a positive
power-of-two modulus.  Thus even a fifth-bit computation, by itself, cannot
decide the required margin condition.

\begin{proposition}[parity does not decide the reduced margin condition]
\label{res:paritysocket}
Under the explicit radius inequalities, parity does not decide membership in
the reduced margin interval, including at a positive power-of-two modulus.
\end{proposition}
\noindent Formalised:
\lword{Totient827ParitySocketCeiling.lean}{73}{reducedMarginSocket_not_decided_by_parity}{parity obstruction}
and
\lword{Totient827ParitySocketCeiling.lean}{110}{powerTwo_reducedMarginSocket_not_decided_by_parity}{power-of-two parity obstruction}.
The result is structural and proves no signed-margin certificate at any
particular scale; it excludes parity as a sufficient standalone criterion.

\paragraph{Prime-adjunction no-go theorem.} One tempting way to seek an
unbounded obstruction is to compare the full target at $H$, $pH$, $qH$, and
$pqH$.  The following affine transport law rules out this route before any
projection is introduced.
For every positive $k$, the actual diagonal satisfies an affine transport law
\[
 D_{kH}=Q_k(H)D_H+Z_k(H),
 \qquad Q_k(H)\in\N,\quad Z_k(H)\in\Z.
\]
Hence integrality at $H$ propagates to every multiple, and the four target hits
at the prime-adjunction diamond are equivalent to the single hit at its root
(\lword{FullTargetPrimeAdjunctionNoGo.lean}{196}{fullTarget_primeAdjunction_diamond_iff_root}{prime-adjunction diamond collapse}).
The corresponding transport of the complete foreign defect has zero diamond
curvature
(\lword{FullTargetPrimeAdjunctionNoGo.lean}{247}{foreignTransportCorrection_primeAdjunction_flat}{foreign-correction flatness}).
Thus explicit-shadow fibres alone cannot create a holonomy obstruction.  A
future positive route would need a separately defined projection together with
a theorem controlling what that projection discards.  This is a proved
obstruction to one strategy, not progress on the unbounded-supply proposition.

\begin{proposition}[scalar-localisation obstruction]\label{res:adelic}
Let $x\in\Q$ and $c\in\Z$.  If $H\mid\operatorname{den}(x)$ and the reduced denominator
of $cx$ divides $H$, then
\[
 \frac{\operatorname{den}(x)}{H}\mid |c|.
\]
\end{proposition}
\noindent Formalised:
\lword{AdelicHeightObstruction.lean}{23}{scalarLocalization_complement_dvd}{complementary-denominator divisibility}
and \lword{AdelicHeightObstruction.lean}{56}{scalarLocalization_integer_eq_mul_num}{scalar-localisation identity}.
Thus a scalar localisation can move a complementary denominator factor into
its coefficient, but cannot erase it.  This rules out a scalar-denominator
shortcut of that form.  It neither proves full-target avoidance nor supplies
any certificates, so it gives no irrationality criterion for $S$.

\begin{proposition}[square-CRT correction suppression]\label{res:squarecrt}
Let $E$ be a finite family indexed by distinct primes $p_i$, with prescribed
anchors $A_i$ and residues $a_i$.  There is a bounded common base $n$ such
that, for each $i\in E$ and each shift $h$ with $p_i\nmid a_i+h$, one has
\[
 n=A_i+p_i(a_i+p_i t)\quad\text{and}\quad
 \ph(n+p_i h-A_i)=(p_i-1)\ph(a_i+p_i t+h)
\]
for a suitable $t\in\N$.  In particular, if every $p_i$ exceeds a fixed
horizon $J$, the choice $a_i=0$ works simultaneously for all $1\le h\le J$.
\end{proposition}
\noindent Formalised:
\lword{SquareCRTCube.lean}{198}{exists_squareCRT_clean_totient_family}{finite clean totient family}
and \lword{SquareCRTCube.lean}{228}{exists_squareCRT_clean_horizon}{clean horizon family}.
The finite mechanism does not force a cube coefficient to be nonzero: there
is a checked clean two-step cube with both displayed coefficients zero
(\lword{SquareCRTCube.lean}{285}{squareCRT_clean_block_can_vanish}{vanishing clean block}) and a
separate checked clean cube with second coefficient $-4$
(\lword{SquareCRTCube.lean}{303}{squareCRT_clean_block_can_be_nonzero}{nonzero clean block}).
Thus the result gives neither an unbounded certificate supply nor an
irrationality criterion for $S$.

\subsection{Finite determinant and rank routes}

\begin{proposition}[signed dyadic-moment algebra]\label{res:signedmoment}
For matrices $M:\iota\times\kappa\to R$ and
$N:\kappa\times\iota\to R$ over a commutative ring, the determinant of
$MN$ is the sum over all maps $p:\iota\to\kappa$ of the determinant of the
corresponding selected-column matrix times the diagonal product
$\prod_i N_{p(i),i}$.  Separately, if a selected integer numerator is odd and
its dyadic exponent is strictly larger than every other selected exponent,
then the common numerator after dyadic denominator clearing is odd and hence
nonzero.
\end{proposition}
\noindent Formalised:
\lword{SignedQMomentObstruction.lean}{28}{det_mul_rectangular}{rectangular determinant expansion},
\lword{SignedQMomentObstruction.lean}{77}{scaled_dyadic_sum_odd}{dominant dyadic parity}, and
\lword{SignedQMomentObstruction.lean}{95}{scaled_dyadic_sum_ne_zero}{dominant dyadic nonvanishing}.
These are finite algebraic inputs to a signed-moment route.  They do not show
that any actual totient Hankel determinant is nonzero, and they give no
irrationality criterion for $S$.

\begin{proposition}[dyadic-kernel rank criterion]\label{res:dyadictotient}
For every fixed level $e$, a separated nonzero evaluation minor of the
canonical dyadic totient-kernel family gives its span dimension
$2^e+1$.  Every even residue channel is exactly a scalar multiple of its
odd-core channel.  Moreover, a compressed-adjoint certificate whose nonzero
multiple is strictly smaller in absolute value than its positive modulus is
inconsistent.
\end{proposition}
\noindent Formalised:
\lword{TotientMahlerDefect.lean}{98}{finrank_canonicalTotientKernel_eq_of_certificate}{separated-minor rank criterion},
\lword{TotientMahlerDefect.lean}{156}{totientKernel_even_residue_reduce}{even-residue reduction}, and
\lword{TotientMahlerDefect.lean}{197}{false_of_compressedAdjointCertificate}{compressed-adjoint contradiction}.
What is missing is an all-level construction of the separated minors, for
example by a CRT--Dirichlet argument.  The proposition does not provide those
certificates for every $e$, and it gives no irrationality criterion for $S$.

\begin{proposition}[residual-gauge obstruction for monomial minors]\label{res:residualgauge}
Let a monomial matrix have a common nonzero residual weight in each column.
The residual weights act by a right diagonal gauge, and therefore preserve
determinant nonvanishing.  If one exponent is $1$, the inverse-phase locked
gauge makes that row identically $1$ while preserving a nonzero determinant
whenever the coefficient-side determinant is nonzero.  If the residual is
allowed to depend on both row and column, it reconstructs a consecutive-power
matrix whose zeroth row is identically $1$.
\end{proposition}
\noindent Formalised:
\lword{ResidualGaugeObstruction.lean}{55}{det_residualMonomialMatrix_ne_zero_iff}{residual-gauge determinant equivalence},
\lword{ResidualGaugeObstruction.lean}{94}{locked_reconstruction_preserves_nonzero_minor}{locked-reconstruction nonvanishing}, and
\lword{ResidualGaugeObstruction.lean}{145}{rowReconstruction_zero_row_one}{reconstructed zero row}.
Thus a residual-blind rank, determinant, or conditioning criterion alone cannot
exclude the target.  The proposition does not exclude a criterion that also
uses an arithmetic coupling identity, and it gives no irrationality criterion
for $S$.

\paragraph{Appendix consequence.}
The first subsection gives exact normal forms for the same unbounded supply
as Theorem~\ref{res:red21}; the second turns those forms into finite
projections and checked instances.  The last two subsections identify
additional arithmetic inputs and prove that several weaker parity,
prime-adjunction, scalar, and residual-blind rank criteria cannot provide
them.  No layer changes the quantifier: a positive proof must still produce
certificates beyond every threshold.

}

\section{Mathematical architecture of the checked arguments}
\label{sec:architecture}

\noindent This section is interpretive: it introduces no additional
mathematical claim.  Its purpose is to identify which parts of the checked
arguments are finite and complete, which parts require a uniform theorem, and
why the chosen coordinates expose that distinction.

\paragraph{Finite arithmetic versus uniform supply.}
The arguments repeatedly separate a finite arithmetic core from a statement
that must hold uniformly.  In the full-support Erd\H{o}s--Borwein theorem, the
finite core consists of a Bertrand--CRT choice, divisor-pair averaging, and
explicit parameter inequalities.  For \#249 the separation is sharper:
eventual binary periodicity predicts an integral tail difference; finite
modular arithmetic can refute that prediction at fixed parameters; and
Proposition~\ref{res:complete} proves that the finite certificate loses no
information there.  The missing ingredient is therefore not a larger finite
verification, but a theorem producing such witnesses without bound.

\paragraph{Tail coordinates and completeness.}
The choice of tail coordinates is essential.  Rather than reason about an
informal repeating digit string, one uses the scaled tails $R_N$ and the fact
that rationality forces suitable period differences to be integers.
Certificate completeness then separates two logically distinct tasks: finite
arithmetic decides non-integrality at a fixed pair $(h,N)$, while
Theorem~\ref{res:red21} requires those pairs beyond every threshold.

\paragraph{Equivalent representations.}
Three further changes of coordinates serve different mathematical purposes. The
Mersenne--Lambert transform places $S$ next to the Erd\H{o}s--Borwein constant
and exposes the nonnegative primitive-conductor weight. The M\"obius-square
identity turns the same number into a signed rapidly convergent series. The
coprime-pair identity turns it into a probability. Formalising all three and
proving that the denominator exclusion transfers between them prevents an
informal change of representation from silently changing the claim.  Reducing
the random pair by its gcd then yields the exact Stern--Brocot cylinder law of
Theorem~\ref{res:sternbrocotcylinders} and the Fibonacci--continuant estimates
of Theorem~\ref{res:sternbrocotruns}.  These are unconditional structure inside
the probability representation, not a substitute for the certificate supply.

\paragraph{Necessary conditions on a rational support.}
The mathematical assertions in this part of the architecture are exactly
those of Theorem~\ref{res:rigidity-main}, with their proofs and auxiliary
coordinates in Appendix~\ref{sec:carry}.  Their role is diagnostic.  The
zero-window bound in Theorem~\ref{res:rigidity-main}\textup{(ii)} is obtained
by transporting the fixed-power divisor estimate
\(\tau(n)^k\le(k^{2^k})^k n\) through the binary tail and the exact carry
recurrence.  Part~\textup{(i)} shows that an infinite rational support cannot
be represented by a bounded carry alphabet.  Part~\textup{(iii)} relates the
odd denominator to reciprocal mass, retaining divergence as an explicit
alternative.  None of these necessary conditions is a contradiction; they
describe the infinite Boolean carry path that a future proof would have to
exclude.

\paragraph{General proof template and residual.}
The same architecture is useful beyond the two named problems: choose a
representation in which rationality imposes a rigid tail law, isolate a finite
predicate that soundly refutes that law, prove the predicate complete at fixed
parameters, and state the remaining uniformity theorem separately. The last
step matters as much as the first three. It prevents a large verified finite
range from being mistaken for the infinite theorem.  Thus the exact
consequence of this section is organisational rather than mathematical:
fixed-parameter completeness has been separated from the unbounded supply
required for \#249 and from the two unresolved \#257 directions.

\section{Formal verification and reproducibility}\label{sec:verify}

\noindent This section records proof authority and reproducibility; it does
not add a mathematical conclusion to either theorem spine.

\paragraph{Proof authority and source identity.}
The Lean~4 proof assistant~\cite{lean} and Mathlib library~\cite{mathlib} check
the named results.  The source, exact informal-to-formal links, and build
instructions are available in the
\href{https://github.com/wcook04/plectis-lean-erdos249-257}{public repository}.
This paper is tied to an
\href{\sourceurl}{archived formal-source revision}; every blue mathematical
caption links directly to the corresponding checked theorem there. The repository
pins the theorem prover and library versions needed to reproduce the build.

\paragraph{Computational and expository boundary.}
Every declaration is elaborated and every proof term is checked by the Lean
kernel.  Finite computations use kernel-checked decision procedures rather
than an external evaluator, and the source contains no proof placeholders or
added axioms.  The manuscript is exposition rather than proof authority;
Appendix~\ref{app:index} explains how each displayed result is connected to
its exact formal statement.

\paragraph{Section consequence.}
The linked declarations and their checked proof terms certify the stated
proved, conditional, finite, and equivalence results at the archived source
revision.  That verification does not change the mathematical status of the
unresolved statements: the open problems collected in
Section~\ref{sec:further} remain open.

\section{Exact frontiers and further questions}\label{sec:further}

The conclusions divide into one open \#249 obligation and a two-way \#257
frontier.  For
Erd\H{o}s \#249, Theorem~\ref{res:farey} excludes a large finite denominator
range and Theorem~\ref{res:red21} gives an exact certificate
characterisation; the missing assertion is the unbounded certificate supply
\eqref{eq:unbounded-certificate-supply}.  For Erd\H{o}s \#257,
Theorem~\ref{res:full} settles full support and
Theorem~\ref{res:support} settles several structured infinite supports, but
the universal statement remains open.  A counterexample would instead be a
rational member of the achievement set $\mathcal A$, with $1/2$ the
distinguished target analysed in Section~\ref{sec:eb}.
The achievement-set topology and the Stern--Brocot cylinder law are exact
structural results on the two sides of this frontier; neither supplies the
missing universal or cofinal assertion.

\long\def\FurtherLambertAppendix{%

\subsection{Additional Lambert, probability, and finite-algebraic structure}
\label{sec:additional-lambert}

\paragraph{Periodic weights.}\label{sec:periodicweights}
For \(L_b(\gamma)=\sum_{a\ge1}\gamma(a)/(b^a-1)\), the checked endpoints are:
\begin{theorem}[eventually-periodic nonnegative weights]
\label{res:periodicnonnegative}
An eventually-periodic nonnegative rational weight with a positive periodic
value gives an irrational series.
\end{theorem}
\noindent Formalised:
\lword{CertificateKernel.lean}{12397}{irrational_weightedErdosSeries_eventuallyPeriodic}{eventually-periodic nonnegative criterion}.
This is contained in a broader theorem of Luca--Tachiya~\cite{lucatachiya}.
\begin{theorem}[periodic signed-weight dichotomy]\label{res:periodicsigned}
For a periodic integer weight, the series is irrational or terminating in
base \(b\).
\end{theorem}
\noindent Formalised:
\lword{CertificateKernel.lean}{13847}{irrational_or_bpow_mul_eq_intCast_intWeightedErdosSeries_periodic}{periodic signed-weight dichotomy}.
\begin{proposition}[residue-blind obstruction]\label{res:periodfourobstruction}
For the period-four weight \(1,0,-1,0\), the divisor coefficient vanishes on
\(n\equiv3\pmod4\).
\end{proposition}
\noindent Formalised:
\lword{CertificateKernel.lean}{13898}{intWeightedCoeff_periodFourSignWeight_eq_zero_of_mod_four_eq_three}{period-four coefficient obstruction}.

\paragraph{Finite algebraic boundaries.}\label{sec:primitiveindex}
Put \(a_n=(\ph*\mu)(n)/n\).  If an integer \(D\) clears these coordinates
through \(N\), the exact obstruction is:
\begin{proposition}[primitive-coordinate finite-jet obstruction]
\label{res:primitiveindex}
\(D\) is divisible by an explicit two-tier primorial \(T_N\); hence no fixed
\(D\) clears every \(a_n\).
\end{proposition}
\noindent Formalised:
\lword{PrimitiveDeterminantLift.lean}{148}{twoTierPrimorial_dvd_of_integral_jet}{two-tier primorial obstruction}.

For \(d\ge1\), define the pointed periodic atom
\[
 \omega_d(N)=\frac{2^{\,N\bmod d}}{2^d-1}.
\]
It satisfies the exact carry law
\[
 2\omega_d(N)-\omega_d(N+1)=\mathbf1_{d\mid N+1}.
\]

\begin{proposition}[finite Mersenne-atom determinant boundary]
\label{res:mersennetailatoms}
Fix \(m\ge1\), conductors \(d_j\ge1\), indices \(n_{ij}\ge0\), and integer
column weights \(w_j\), and set
\[
 W_{ij}=w_j\omega_{d_j}(n_{ij}),
 \qquad U_{ij}=2^{\,n_{ij}\bmod d_j}.
\]
Then
\[
 \left(\prod_{j=1}^m(2^{d_j}-1)\right)\det W
 =\left(\prod_{j=1}^m w_j\right)\det U\in\Z.
\]
If the integer on the right is nonzero, then the absolute value of the
cleared determinant on the left is at least \(1\).
\end{proposition}
\noindent Formalised:
\lword{MersenneTailAtoms.lean}{28}{mersenneTailAtom_step}{Mersenne-atom carry law},
\lword{MersenneTailAtoms.lean}{158}{weightedMersenneAtom_denominatorProduct_mul_det}{determinant-clearing identity},
and
\lword{MersenneTailAtoms.lean}{181}{one_le_abs_weightedMersenneAtom_cleared_det}{cleared-determinant lower bound}.
Thus finite Mersenne-atom determinants recover exactly their explicit column
denominators; the clearing creates no additional height surplus.
These are scoped finite-denominator obstructions, not irrationality criteria
for \(S\).

\begin{proposition}[the M\"obius-square lens]\label{res:lens}
\[
S \;=\; \tfrac12 \;+\; \sum_{d\ge1}\frac{\mu(d)}{(2^d-1)^2}.
\]
\end{proposition}
\noindent Formalised: \lword{CertificateKernel.lean}{18126}{totient_series_eq_half_add_moebius_mersenne_square}{M\"obius-square identity}.
Writing $T=\sum_{d\ge1}\mu(d)/(2^d-1)^2$, irrationality of $S$ is equivalent to
irrationality of $T$, since $S=\tfrac12+T$. The factor $1/(2^d-1)^2$ is the
probability that $d$ divides two independent fair-coin waiting times
(Section~\ref{sec:geom}). The signed constant, the coprime-pair mass, and the
tail differences of Section~\ref{sec:reduction} are three readings of the same
quantity.

\paragraph{Squared Lambert calculus.}\label{sec:squared}
For \(L_2(f)=\sum_{d\ge1}f(d)/(2^d-1)^2\), the denominator is the probability
that \(d\) divides two independent geometric waiting times.  The formal
transfer theorem,
\lword{GcdMomentCalculus.lean}{105}{tsum_lambert_linear_weight_sq_pure}{squared Lambert transfer},
therefore gives a divisor calculus for their gcd.
\begin{proposition}[squared Lambert ladder]\label{res:gcdmoment}
At \(r=1/2\), the three rows \(f=\mu,\mathbf 1,\ph\) yield respectively
\[
 L_2(\mu)=S-\tfrac12,\qquad
 L_2(\mathbf 1)=\sum_{n\ge1}\frac{\sigma(n)-\tau(n)}{2^n},\qquad
 L_2(\ph)=\sum_{n\ge1}\frac{P(n)-n}{2^n},
\]
where \(P\) is Pillai's gcd-sum function.
\end{proposition}
\noindent Formalised:
\lword{GcdMomentCalculus.lean}{235}{tsum_totient_div_mersenne_sq_eq_gcd_moment_series}{totient gcd-moment identity}.
Postelmans--Van Assche imply irrationality of the middle row
~\cite{postelmansvanassche}; this does not transfer to the M\"obius row, which
is exactly the open constant \(S-\tfrac12\).

\paragraph{Reduced slopes and Stern--Brocot cylinders.}
The fair-coin pair has a particularly rigid distribution after division by
its gcd.  For positive coprime \(a,b\), put
\[
  \nu(a,b)=\frac1{2^{a+b}-1},
  \qquad
  M(a,b)=\frac1{(2^a-1)(2^b-1)}.
\]
The first quantity is the probability that the reduced direction is
\((a,b)\); the second is the mass of the Stern--Brocot cylinder rooted at
\((a,b)\).  Let \(M_d(a,b)\) be the total stop mass in the first \(d\)
generations below that root, where a node \((u,v)\) has stop mass
\((2^{u+v}-1)^{-1}\) and children \((u+v,v)\), \((u,u+v)\).

\begin{theorem}[reduced-direction and cylinder law]
\label{res:sternbrocotcylinders}
The reduced directions form a probability distribution,
\[
  \sum_{\substack{a,b\ge1\\(a,b)=1}}\nu(a,b)=1.
\]
At every positive node,
\[
  M(a,b)=\frac1{2^{a+b}-1}+M(a+b,b)+M(a,a+b),
\]
and, uniformly in \(a,b\) and \(d\),
\[
  0\le M(a,b)-M_d(a,b)
  \le \left(\frac23\right)^d M(a,b).
\]
Consequently \(M_d(a,b)\to M(a,b)\); at the root \(M(1,1)=1\).
\end{theorem}
\noindent Formalised:
\lword{GcdMomentCalculus.lean}{349}{tsum_pos_coprime_inv_mersenne_eq_one}{reduced-direction mass one},
\lword{GcdMomentCalculus.lean}{474}{cylinderMass_split}{cylinder mass recursion},
\lword{GcdMomentCalculus.lean}{525}{sternBrocotDepthMass_error}{depth-$d$ remainder bound}, and
\lword{GcdMomentCalculus.lean}{560}{tendsto_sternBrocotDepthMass}{cylinder convergence}.
The recursion is the elementary rational identity obtained after setting
\(A=2^a\), \(B=2^b\); the quantitative point is that the two children retain
at most two thirds of their parent's mass.

\paragraph{Alternating-run stability.}
Write an alternating Stern--Brocot word as \(r\) nonempty runs of lengths
\(n_i=1+e_i\), with \(e_i\ge0\).  Starting with \(P_\varnothing=(1,1)\), define
\[
  P_{(n_1,\ldots,n_r)}=(n_1A+B,A)
  \quad\text{when}\quad P_{(n_2,\ldots,n_r)}=(A,B),
  \qquad H(n_1,\ldots,n_r)=P_1+P_2.
\]
Thus \(H\) is the arithmetic height of the primitive pair reached by the
word.  Let \(F_0=0,F_1=1\).

\begin{theorem}[Fibonacci--continuant stability]
\label{res:sternbrocotruns}
Every \(r\)-run word satisfies
\[
 H(1+e_1,\ldots,1+e_r)
 \ge F_{r+3}+F_{r+1}\sum_{i=1}^r e_i,
\]
and \(H(1,\ldots,1)=F_{r+3}\).  A defect confined to position
\(i\in\{0,\ldots,r-1\}\) has the exact size
\[
 H(\underbrace{1,\ldots,1}_{i},1+e,
   \underbrace{1,\ldots,1}_{r-i-1})
 =F_{r+3}+F_{i+2}F_{r-i+1}e.
\]
More generally, the checked recurrence gives the full positive multiaffine
continuant defect.  The two orientations with one nonempty run have total
mass \(1/2\), while the natural denominator exponent on the all-unit spine is
\[
  \sum_{j=0}^{r-1}F_{j+2}=F_{r+3}-2.
\]
\end{theorem}
\noindent Formalised:
\lword{SternBrocotRunGeometry.lean}{185}{runHeight_fib_lower}{Fibonacci height lower bound},
\lword{SternBrocotRunGeometry.lean}{263}{runHeight_defectRunLengths}{multiaffine run-height formula},
\lword{SternBrocotRunGeometry.lean}{343}{runHeight_defect_fib_sum_lower}{aggregate defect lower bound},
\lword{SternBrocotRunGeometry.lean}{413}{runHeight_oneSiteDefect}{one-site defect identity},
\lword{SternBrocotRunGeometry.lean}{467}{oneRunMass_eq_half}{one-run mass}, and
\lword{SternBrocotRunGeometry.lean}{491}{naturalRunDenominatorExponent_add_two}{natural denominator exponent}.

The Lambert and M\"obius identities used above are classical
~\cite{apostol,merca2017,mercaschmidt}, and the mediant language goes back to
Farey~\cite{farey1816}.  Postelmans--Van Assche supply the cited irrationality
of the constant-weight squared-Lambert row, not these cylinder or run
statements.  We claim the displayed Lean-checked identities and estimates,
not priority for their assembled package.  In particular, the exact equality
of the Fibonacci height and denominator scales is a boundary, not an
irrationality argument: no theorem here proves that a run tail survives
denominator clearing.
}

\subsection{Two concrete scalar criteria}

Two scalar conditions give concrete sufficient routes toward the two open
directions.  Both are exact arithmetic statements; neither is proved at the
unbounded range needed for its final conclusion.

For \#249, let $m_t$ be the canonical adjacent-suffix depth and let $d_t$ be
the corresponding residue in $[0,2^{m_t})$.  Define
\[
  \sigma_t=\min\bigl(d_t-2^{m_t-5},
    (2^{m_t}-2^{m_t-5})-d_t\bigr)\in\Z.
\]

\begin{proposition}[canonical strict-jump slack]\label{res:slackfrontier}
The canonical residue lies in the central band
$[2^{m_t-5},2^{m_t}-2^{m_t-5}]$ if and only if $\sigma_t\ge0$.
If nonnegative slack occurs at arbitrarily large strict jumps
$M_t<M_{t+1}$, then $S$ is irrational.
\end{proposition}
\noindent Formalised:
\lword{DiagonalFreshLossBridge.lean}{2469}{canonicalAdjacentSuffixCentral_iff_slack_nonneg}{canonical central-band equivalence},
\lword{DiagonalFreshLossBridge.lean}{2559}{canonicalAdjacentSuffixJumpCentralSupply_iff_slackSupply}{strict-jump slack equivalence},
and
\lword{DiagonalFreshLossBridge.lean}{2731}{irrational_totientSeries_of_canonicalAdjacentSuffixJumpSlackSupply}{cofinal-slack irrationality criterion}.
The proposition is an equivalence followed by a sufficient criterion.  It
does not prove that $\sigma_t$ is nonnegative cofinally.

For \#257, write $r_n$ for the exact greedy residual of the target $1/2$ and
put
\[
  u_n=4^n(2r_n-2^{-n}).
\]
The second-channel condition is the rational inequality
\[
  \frac16+\frac{37}{56}2^{-n}\le \left|u_n-\frac13\right|.
\]

\begin{proposition}[rational second-channel reduction]\label{res:secondfrontier}
The second-channel condition is decidable over $\Q$ and holds for
$1\le n\le6$.  If it holds for every $n\ge7$, then
$1/2$ belongs to the Mersenne achievement set $\mathcal A$.
\end{proposition}
\noindent Formalised:
\lword{GreedyAchievementSet.lean}{2695}{HalfSecondChannelSeparatedRat}{second-channel predicate},
\lword{GreedyAchievementSet.lean}{2723}{halfSecondChannelSeparatedRat_of_pos_le_six}{finite second-channel verification},
and
\lword{GreedyAchievementSet.lean}{2901}{half_mem_mersenneAchievementSet_of_secondChannelSeparationRat_from_seven}{from-seven membership criterion}.
The membership conclusion is itself exact:
\[
  1/2\in\mathcal A
  \quad\Longleftrightarrow\quad
  \text{the canonical greedy orbit omits infinitely many exponents}.
\]
This is formalised at
\lword{GreedyAchievementSet.lean}{2514}{half_mem_mersenneAchievementSet_iff_greedySkippedSupport_infinite}{half-membership--infinite-skip equivalence}.
The forward implication uses Theorem~\ref{res:full}; the reverse implication
uses the absorbing fatal-state
criterion.  Thus the second-channel hypothesis is one sufficient route to the
infinite-skip condition.  It does not prove that condition, and membership alone
is not the universal \#257 theorem.

\subsection{The unresolved statements}

\begin{problem}[Erd\H{o}s \#249]
Prove that
\[
S=\sum_{n\ge1}\frac{\ph(n)}{2^n}
\]
is irrational.
\end{problem}

By Theorem~\ref{res:red21}, the preceding problem is equivalent to the
following exact quantifier statement.

\begin{problem}[unbounded certificate supply]
Prove~\eqref{eq:unbounded-certificate-supply}.
The lcm-diagonal supply in Corollary~\ref{res:diag} is sufficient, but the
finite scales in Example~\ref{res:deposits} do not establish it.
\end{problem}

\begin{problem}[universal Erd\H{o}s \#257]
Prove that
\[
  \sum_{n\in A}\frac{1}{2^n-1}
\]
is irrational for every infinite set $A\subseteq\N$, rather than only for the
support families in Theorem~\ref{res:support}.
\end{problem}

\begin{problem}[the half-value counterexample route]
Decide whether $1/2\in\mathcal A$.  By
Theorem~\ref{res:halfmembership}, this is equivalent to the canonical greedy
orbit omitting infinitely many exponents.  A positive answer would give an
infinite support of rational value $1/2$ and refute the universal \#257
statement; a negative answer would close this distinguished counterexample
route only.
\end{problem}

The scalar criteria above isolate two concrete subproblems: cofinal
nonnegative slack would settle \#249, while the second-channel inequality
would settle the half-value membership question.  Neither assertion follows
from its checked finite range.

Rationality on the \#257 side would also impose two complementary forms of
rigidity: divisor-coverage zero windows have length $o(\log N)$, whereas an
infinite support with rational series value forces the associated positive
carry states to be unbounded.  These constraints do not yet exclude a
rational infinite-support value.

Appendix~\ref{app:lambert-complements} collects further weighted-Lambert,
M\"obius-square, and finite-algebraic identities.  They provide additional
coordinates and obstructions, but no further closure of the three statements
above.  Across both problems the missing step has the same logical form: a
finite obstruction is exact at fixed parameters, while the theorem requires
that obstruction uniformly or without bound.

\subsection*{Statements and declarations}

\paragraph{Artefact and data availability.} The
\href{\sourceurl}{archived formal-source revision} contains the Lean sources,
pinned toolchain, library manifest, and generated certificate data used in the
verification. Repository metadata and this manuscript provide navigation
rather than proof authority; Section~\ref{sec:verify} gives the verification
route.

\paragraph{Use of AI assistance.} Large-language-model assistants were used
during development for proof drafting, refactoring, and prose editing. Every
registered formal claim is linked to a declaration in the archived source
revision. Lean checks each proof term against the pinned library; the project
contains no proof placeholders or project-defined axioms.

\paragraph{Funding and competing interests.} This work received no external
funding. The author declares no competing interests.

\paragraph{Acknowledgements.} The problem numbering and status follow the
Erd\H{o}s Problems catalogue maintained by Thomas Bloom~\cite{erdosproblems}.

\appendix
\section{Technical reductions for the half-value branch}
\label{app:half-technical}

This appendix records the local arithmetic used to reach
Theorems~\ref{res:halfmembership} and~\ref{res:fatalright}.  Its statements
fall into four classes: exact coordinate identities, conditional compactness
criteria, bounded finite certificates, and counterexamples to tempting local
inductions.  None changes the two open conclusions at the end of
Section~\ref{sec:eb}.  The appendix is included so that the hypotheses in the
main reduction can be audited without turning the main theorem section into a
chronology of the formal development.

The order is by mathematical use.  Appendix~A.1 localizes one possible unsafe
skip; Appendix~A.2 records conditional structured-support criteria;
Appendix~A.3 collects the finite-window and compactness implications;
Appendix~A.4 develops the equivalent final-skip coordinates; and
Appendix~A.5 records verified finite instances, exact coordinate
translations, counterexamples, and closed strategy branches.  This is a
dependency map, not a second route to a conclusion.

Throughout this appendix, a \emph{seam} is an adjacent pair of finite greedy
words at a branch boundary, a \emph{rewind} is the corresponding inverse
finite recurrence, and the \emph{transition carry} is the integer coordinate
appearing in the exact residual identity.  These terms denote mathematical
objects, not stages of a proof search.

\HalfValueTechnicalAppendix

\paragraph{Appendix consequence.}
The compactness and final-skip subsections give exact implications: a cofinal
window, seam, or terminal approximation would produce \(1/2\in\mathcal A\),
whereas an eventual right tail is exactly nonmembership.  The other
subsections localize possible inputs, verify bounded instances, and close
stronger but invalid local inferences.  None produces the required cofinal
input, so the residual remains precisely the two obligations in
Theorem~\ref{res:fatalright}.

\section{Auxiliary rationality criteria from binary carry systems}\label{sec:carry}

Five further theorem families provide additional criteria. They prove no new
irrationality. Instead they pin down, in machine-checked form, what a rational
value of the series in question would have to look like: as an integer carry
orbit, as a M\"obius-inverted support coordinate, as a residue-wrap cycle
modulo the odd part of the denominator, and as a point of a greedy subset-sum
geometry. Everything in this section is an equivalence or an unconditional
constraint on the rational side; nothing excludes a rational value of either
open series, and each family states that boundary explicitly. The prior-art
boundary is theorem-family-specific. The integral carry recurrence has a
direct antecedent in Wang--Grau Ribas~\cite{wanggrauribas}; the strict-tail
geometry is recorded by Kova\v{c}--Tao~\cite{kovactao}; M\"obius inversion,
repetend algebra, and divisor averaging are classical (see, for example,
Apostol~\cite{apostol}). We make no claim of mathematical priority for the
converse/rigidity, certificate-normal-form, or coupled reciprocal-mass/collision
statements collected here.

The reading order is: a generic tail-orbit rationality normal form in
Appendix~B.1; its Boolean--M\"obius specialization in Appendix~B.2;
divisor-coverage constraints in Appendix~B.3; denominator-wrap, reciprocal
mass, and unbounded-state constraints in Appendix~B.4; and the exact greedy
geometry with one-sided finite death certificates in Appendix~B.5.

\subsection{Tail-orbit rigidity: a rationality criterion}

For a coefficient sequence $c:\N\to\N$ with $c(n)\le n$, write
\[
X_c=\sum_{n\ge1}\frac{c(n)}{2^n},
\qquad
T_c(N)=\sum_{j\ge1}\frac{c(N+j)}{2^j},
\]
the series and its scaled tail after $N$ binary digits. Call an integer
sequence $u$ a \emph{tempered orbit} for $c$ with multiplier $v\ge1$ if
\[
u(N+1)=2\,u(N)-v\,c(N+1)\ \text{ for all } N,
\qquad\text{and}\qquad
u(N)/2^N\to0 .
\]
The recurrence is exact binary long division against the coefficient stream;
the boundary condition forbids the homogeneous solution. Wang and Grau Ribas
use the integral carry recurrence forced by rationality for the weighted binary
special case $c(n)=n d_n$~\cite{wanggrauribas}. The theorem below extends that
forward mechanism to arbitrary nonnegative $c(n)\le n$; its explicit converse
and subexponential-orbit rigidity are the local additions, with no priority
claim attached. Wang--Grau Ribas's positive-density theorem and its Erd\H{o}s
\#260 corollary are not used or formalised here.

\begin{theorem}[tail-orbit rigidity]\label{res:rigidity}
Let $c(n)\le n$ for all $n$. Every tempered orbit is the scaled tail,
$u(N)=v\,T_c(N)$ for all $N$; and $X_c$ is rational if and only if a tempered
orbit exists for some multiplier $v\ge1$.
\end{theorem}
\noindent Formalised:
\lword{GenericTailOrbitRigidity.lean}{338}{temperedBinaryOrbit_eq_scaledTail}{the scaled-tail orbit identity}
(rigidity) and
\lword{GenericTailOrbitRigidity.lean}{425}{binaryCoeffSeries_rational_iff_exists_temperedBinaryOrbit}{the rationality--orbit equivalence}
(the criterion), restated in Mathlib's irrationality vocabulary at
\lword{GenericTailOrbitRigidity.lean}{434}{not_irrational_binaryCoeffSeries_iff_exists_temperedBinaryOrbit}{the non-irrationality--orbit equivalence}.
The engine is one observation: a real orbit with $d(N+1)=2\,d(N)$ and
$d(N)=o(2^N)$ vanishes identically
(\lword{GenericTailOrbitRigidity.lean}{314}{doublingOrbit_eq_zero_of_tempered}{zero-tail rigidity}).

The tempered boundary is essential: the
homogeneous parasite $N\mapsto2^N$ can be added to any orbit without
disturbing the recurrence, so positivity of an orbit certifies nothing by
itself, and positivity is nowhere advertised as an equivalent criterion. The
linear-growth hypothesis covers both constants introduced in
Section~\ref{sec:intro}:
$\ph(n)\le n$ for the \#249 coefficients, and $f_A(n)\le\tau(n)\le n$ for the
\#257 support coefficients $f_A(n)=\#\{a\in A:a\mid n\}$. The results below
instantiate the criterion on supports; the totient instantiation is not taken
here.

\begin{proposition}[finite-level carry anti-compression]\label{res:carryrank}
Suppose a tempered integral carry orbit for the totient sequence has positive
multiplier.  If the canonical level-$e$ totient kernel family is linearly
independent, then the span of its canonical carry sections has dimension at
least $2^e-1$.  Consequently, an all-level independence hypothesis together
with rationality of the totient coefficient series would force one tempered
carry orbit with unbounded finite-level section ranks.
\end{proposition}
\noindent Formalised:
\lword{TotientCarryKernelRigidity.lean}{211}{finrank_canonicalCarryKernel_ge_of_linearIndependent}{the carry-kernel rank lower bound}
and
\lword{TotientCarryKernelRigidity.lean}{261}{not_irrational_totientSeries_implies_unbounded_carryRank}{the unbounded-rank consequence of rationality}.
The all-level independence hypothesis is not proved here.  Thus the
proposition records an exact conditional anti-compression consequence, not an
irrationality criterion for $S$.

\subsection{Boolean--M\"obius carry coordinates (the \#257 direction)}

At base $2$ the support series of Section~\ref{sec:eb} is itself a binary
coefficient series, $\sum_{n\in A}1/(2^n-1)=X_{f_A}$
(\lword{BooleanMobiusCarry.lean}{376}{erdosSupportSeries_two_eq_binaryCoeffSeries}{the support-series/binary-series identity}).
Two exact coordinate changes then turn a hypothetical rational value into a
support-free object. The divisor transform and M\"obius inversion below are
classical Lambert-series tools~\cite{apostol,merca2017,mercaschmidt}; the
contribution here is their Lean-checked composition with the tempered-orbit
criterion.

\begin{proposition}[M\"obius-sign support no-go]\label{res:mobiussignnogo}
Let
\[
 A_-=\{d\ge2:\mu(d)=-1\},
 \qquad A_+=\{d\ge2:\mu(d)=1\}.
\]
Then
\[
 \sum_{d\in A_-}\frac1{2^d-1}
 =\frac12+\sum_{d\in A_+}\frac1{2^d-1}
 \ge\frac12+\frac1{63}>\frac12.
\]
In particular, the signed identity \(L(\mu)=1/2\) cannot be converted into a
Boolean support of value \(1/2\) merely by selecting the negative M\"obius
indices.
\end{proposition}
\noindent Formalised:
\lword{MobiusSignSupportNoGo.lean}{111}{tsum_negativeMobius_eq_half_add_positiveMobiusTail}{the negative-M\"obius decomposition}
and
\lword{MobiusSignSupportNoGo.lean}{150}{half_add_one_div_sixty_three_le_tsum_negativeMobius}{the strict lower bound}.
The term \(d=6\) supplies the displayed \(1/63\).  This closes one
sign-truncation route only; it does not exclude another infinite Boolean
support with value \(1/2\).

\begin{proposition}[Boolean M\"obius inversion, both directions]\label{res:boolmob}
On positive integers, $f_A=\mathbf 1_A*\zeta$ and $\mu*f_A=\mathbf 1_A$, where
$\mathbf 1_A$ is the indicator of $A$. Conversely, every integer-valued
arithmetic function $f$ whose M\"obius transform is Boolean,
$(\mu*f)(n)\in\{0,1\}$ for all $n\ge1$, satisfies $f=f_B$ for the support
$B=\{n\ge1:(\mu*f)(n)=1\}$ selected by that transform.
\end{proposition}
\noindent Formalised:
\lword{BooleanMobiusCarry.lean}{76}{positiveSupportBitAF_mul_zeta}{the positive-support zeta identity},
\lword{BooleanMobiusCarry.lean}{94}{moebius_mul_supportCoeffAF}{the M\"obius inversion identity}, and, with no
search or finiteness hypothesis, the converse
\lword{BooleanMobiusCarry.lean}{180}{eq_supportCoeffAF_booleanMobiusSupport_of_boolean}{the Boolean reconstruction theorem};
the selected support is
\lword{BooleanMobiusCarry.lean}{139}{booleanMobiusSupport}{the Boolean M\"obius support}.

On the carry side, Theorem~\ref{res:rigidity} applies at $c=f_A$: the support
series is rational exactly when a tempered carry orbit for $f_A$ exists
(\lword{BooleanMobiusCarry.lean}{383}{erdosSupportSeries_rational_iff_exists_temperedCarry}{the rationality--carry equivalence});
a displayed fraction $p/q$ corresponds to an orbit $U$ with multiplier $q$ and
initial value $U(0)=p$
(\lword{BooleanMobiusCarry.lean}{466}{support_fraction_iff_exists_temperedCarry}{the support-fraction carry equivalence});
and exact division recovers the coefficient from the orbit,
$\bigl(2U(N)-U(N+1)\bigr)/q=f_A(N+1)$
(\lword{BooleanMobiusCarry.lean}{680}{temperedCarry_quotient_eq_supportCoeff}{carry-quotient recovery}).
The elementary divisor-pair bound $\tau(n)\le2\lfloor\sqrt n\rfloor$
(\lword{BooleanMobiusCarry.lean}{208}{card_divisors_le_two_mul_sqrt}{the divisor-count bound}) confines
the tails to a square-root strip, $T_{f_A}(N)\le2\sqrt N+4$
(\lword{BooleanMobiusCarry.lean}{289}{binaryCoeffTail_supportCoeff_le_two_sqrt_add_four}{the coefficient-tail bound}),
hence every orbit state to $0<U(N)\le q\,(2\sqrt N+4)$
(\lword{BooleanMobiusCarry.lean}{690}{temperedCarry_pos_of_exists_pos_mem}{carry positivity},
\lword{BooleanMobiusCarry.lean}{705}{temperedCarry_le_denominator_mul_two_sqrt_add_four}{the carry upper bound});
and the strip is strong enough to re-derive the tempered boundary, so the
analytic side condition can be traded for one inequality. The two coordinate
changes compose.

\begin{theorem}[certificate-level equivalence]\label{res:carrycert}
Fix an integer $p$ and an integer $q\ge1$. There is a nonempty support $A$ of
positive integers with $\sum_{n\in A}1/(2^n-1)=p/q$ if and only if there is an
integer sequence $U$ with $U(0)=p$ such that: every $U(N)$ is positive;
$U(N)\le q\,(2\sqrt N+4)$; $q$ divides $2U(N)-U(N+1)$ for every $N$; and the
M\"obius transform of the quotient sequence $\bigl(2U(N)-U(N+1)\bigr)/q$ is
Boolean. The support is not guessed: it is reconstructed from the certificate
by Proposition~\ref{res:boolmob}.
\end{theorem}
\noindent Formalised:
\lword{BooleanMobiusCarry.lean}{948}{exists_normalized_support_fraction_iff_exists_booleanMobiusCarry}{the normalized support-fraction certificate equivalence},
with the certificate packaged as
\lword{BooleanMobiusCarry.lean}{775}{BooleanMobiusCarryCertificate}{the Boolean carry certificate} and the
quotient normalized at index zero
(\lword{BooleanMobiusCarry.lean}{766}{carryQuotient}{the carry quotient}).

\begin{example}[the support $\{2,3\}$]
The support $\{2,3\}$ has value
$\tfrac13+\tfrac17=\tfrac{10}{21}$
(\lword{BooleanMobiusCarry.lean}{1054}{erdosSupportSeries_support23_eq_ten_div_twenty_one}{the two-point series value}),
its orbit is the pure period-six cycle $10,20,19,17,13,26$ with multiplier
$21$ (\lword{BooleanMobiusCarry.lean}{1103}{carryOrbit23_isTempered}{the tempered two-point orbit}), and the
M\"obius transform of the carry quotient recovers exactly $\{2,3\}$
(\lword{BooleanMobiusCarry.lean}{1140}{mobius_carryOrbit23_recovers_support}{support recovery}).
\end{example}

These are coordinates, not by themselves a strategy: nothing here excludes
infinite Boolean carry paths, and no finite search is promoted to a
completeness argument.

\subsection{Sublogarithmic divisor coverage}

The coefficient $f_A(n)$ counts the elements of $A$ dividing $n$. Say that a
zero window of length $h$ starts after $c+N$ when
$f_A(c+N+j+1)=0$ for every $0\le j<h$. Thus no integer in that interval is
divisible by an element of the support.

\begin{theorem}[sublogarithmic zero windows]\label{res:sublog}
Let $A$ contain a positive integer and suppose
\[
  \sum_{a\in A}\frac1{2^a-1}=\frac{p}{2^c v},
  \qquad p\in\Z,\quad c\in\N,\quad v\ge1.
\]
For every $\varepsilon>0$ there is a constant $B\ge0$, depending only on
$\varepsilon,c,$ and $v$, such that every zero window after $c+N$ satisfies
\[
  h\le \varepsilon\log_2(N+1)+B,
  \qquad \log_2 x:=\frac{\log x}{\log 2}.
\]
The bound is uniform in $A,p,N,$ and $h$.
\end{theorem}
\noindent Formalised:
\lword{SublogDivisorCoverage.lean}{392}{supportCoeffZeroWindow_length_le_eps_logb_add}{sublogarithmic zero-window bound}.

The proof first establishes, for each integer $k\ge2$, the explicit divisor
estimate
\[
  \tau(n)^k\le \bigl(k^{2^k}\bigr)^k n.
\]
It transports the resulting $n^{1/k}$ bound through the binary tail and then
uses the exact carry recurrence across a zero window. Raising the resulting
inequality to the $k$th power absorbs the occurrence of $h$ on the right and
leaves a coefficient $1/(k-1)$ in front of the binary logarithm. Choosing $k$
after $\varepsilon$ gives the theorem. This controls gaps in divisor coverage;
it does not assert a subpower bound for the carry state itself.

\subsection{Residue wraps, reciprocal mass, and unbounded carries}

Suppose now the support series equals $p/(2^c v)$ with $v$ odd and $A$
containing a positive element. Rationality makes the scaled tails integers:
$u(N)=v\,T_{f_A}(c+N)$ is a positive integer, satisfies the carry recurrence
$u(N+1)+v\,f_A(c+N+1)=2\,u(N)$, and reduces modulo $v$ to the doubling
residues $r_N=p\,2^N\bmod v$
(\lword{RationalSupportCarrySkeleton.lean}{668}{exists_shifted_odd_tail_nat_states_of_support_fraction}{shifted natural-state existence}).
Doubling a least residue either stays under $v$ or wraps past it exactly once,
so each step emits a wrap digit $k_N=\lfloor2r_N/v\rfloor\in\{0,1\}$, and over
any cycle length $h$ with $2^h\equiv1\pmod v$ the repetend identity holds:
\[
\sum_{N<h}r_N \;=\; v\cdot w,
\qquad
w=\sum_{N<h}k_N,
\]
with $w\ge1$ when $p$ is coprime to $v>1$. Formalised:
\lword{RationalSupportCarrySkeleton.lean}{171}{sum_doublingResidue_eq_mul_wrapCount}{wrap-count identity}
and \lword{RationalSupportCarrySkeleton.lean}{181}{doublingWrapCount_pos}{positive wrap count};
residues, digits, and counts are
\lword{RationalSupportCarrySkeleton.lean}{35}{doublingResidue}{doubling residue},
\lword{RationalSupportCarrySkeleton.lean}{39}{doublingWrapDigit}{wrap digit},
\lword{RationalSupportCarrySkeleton.lean}{43}{doublingWrapCount}{wrap count}.

\begin{proposition}[one-wrap classification]\label{res:onewrap}
Let $v>1$ be odd, $p$ coprime to $v$, and $h\ge1$ with $2^h\equiv1\pmod v$.
If the cycle of $p$ wraps exactly once, then $v=2^h-1$ and the starting
residue is a power of two: $r_0=2^a$ for some $a<h$.
\end{proposition}
\noindent The one-wrap cycles are thus exactly the Mersenne repetends
$2^a/(2^h-1)$. Formalised:
\lword{RationalSupportCarrySkeleton.lean}{415}{one_doublingWrap_classification}{one-wrap classification},
from the generic closed-cycle form
\lword{RationalSupportCarrySkeleton.lean}{353}{one_wrap_cycle_classification}{one-wrap cycle classification},
and instantiated at the true order $h=\operatorname{ord}_v(2)$
(\lword{RationalSupportCarrySkeleton.lean}{478}{one_oddOrder_doublingWrap_classification}{odd-order one-wrap classification},
with \lword{RationalSupportCarrySkeleton.lean}{436}{oddDoublingOrder}{odd doubling order}). The
classification is algebraic; a finite check of $446$ coprime starting residues
across twelve modulus/order rows is retained as kernel-checked validation, not
as the proof
(\lword{RationalSupportCarrySkeleton.lean}{2450}{orderWrapValidatedStartCount_eq_446}{446-start count},
\lword{RationalSupportCarrySkeleton.lean}{2454}{orderWrapValidation446_passes}{446-start validation},
\lword{RationalSupportCarrySkeleton.lean}{2460}{orderWrap_minima_table_passes}{minima-table validation}).

The analytic bridge is Ces\`aro averaging. Write $\rho(A)=\sum_{a\in A}1/a$
for the reciprocal mass
(\lword{RationalSupportCarrySkeleton.lean}{790}{reciprocalMass}{reciprocal mass}). When the
reciprocal family is summable, the mean of $f_A(1),\dots,f_A(N)$ converges to
$\rho(A)$ --- the density of the multiples of $a$, summed over the support ---
and the mean of the scaled tails has the same limit
(\lword{RationalSupportCarrySkeleton.lean}{945}{tendsto_supportCoeff_mean_reciprocalMass}{divisor-mean limit},
\lword{RationalSupportCarrySkeleton.lean}{1036}{tendsto_binaryCoeffTail_supportCoeff_mean_reciprocalMass}{coefficient-tail mean limit}).
The divergent case is carried as an explicit disjunction, not through a
default value assigned to a divergent formal sum
(\lword{RationalSupportCarrySkeleton.lean}{796}{ReciprocalMassDivergesOrAtLeast}{reciprocal-mass dichotomy}).
Each tail state splits exactly as $u(N)=v\,e_N+r_N$ with integral excess
$e_N=\lfloor u(N)/v\rfloor\ge0$
(\lword{RationalSupportCarrySkeleton.lean}{501}{oddTailExcess}{odd-tail excess}).

\begin{proposition}[order-wrap lower bound, with exact excess mean]\label{res:orderwrap}
Let $A$ and $p/(2^c v)$ be as above with $v>1$ odd, and let $w$ be the wrap
count of $p$ over one cycle of length $h=\operatorname{ord}_v(2)$. Then
$\sum_{a\in A}1/a$ diverges or $\rho(A)\ge w/h$; if moreover $p$ is coprime to
$v$, then $w\ge1$, so $\rho(A)\ge1/\operatorname{ord}_v(2)$. In the summable
case the bound is the periodic part of an exact identity: the Ces\`aro mean of
the excess $e_N$ converges automatically, and
\[
\rho(A)\;=\;\frac{w}{h}\;+\;\lim_{N\to\infty}\frac1N\sum_{n<N}e_n .
\]
\end{proposition}
\noindent Formalised:
\lword{RationalSupportCarrySkeleton.lean}{1454}{support_fraction_oddOrder_wrapRatio_le_reciprocalMass}{wrap-ratio mass bound},
\lword{RationalSupportCarrySkeleton.lean}{1478}{one_div_oddOrder_le_reciprocalMass_of_support_fraction}{odd-order mass lower bound},
the form retaining the divergent alternative
\lloc{RationalSupportCarrySkeleton.lean}{1582},
and the excess mean
\lword{RationalSupportCarrySkeleton.lean}{1334}{tendsto_oddTailExcess_mean}{odd-tail excess mean limit}
(fraction-facing shifted form
\lword{RationalSupportCarrySkeleton.lean}{1577}{exists_shifted_oddTailExcess_mean_of_support_fraction}{shifted excess-mean realisation},
named-limit form
\lword{RationalSupportCarrySkeleton.lean}{1353}{reciprocalMass_eq_wrapRatio_add_oddTailExcessMean}{exact mass decomposition}).
A rational value thus ties the odd part of its denominator to the sparsity of
its support: a summable support with small $\rho(A)$ can only display
denominators whose odd part has small multiplicative order of $2$.

\begin{corollary}[collision strengthening at common multiples]\label{res:collision}
In the summable case, let $F\subseteq A$ be finite and let $L$ be a common
multiple of $F$. Then
\[
\rho(A)\;\ge\;\frac{w}{h}\;+\;\frac{\lceil(|F|-1)/2\rceil}{L}.
\]
In particular, an infinite support whose value is a dyadic rational $p/2^c$
has $\sum_{a\in A}1/a$ divergent or $\rho(A)>1$.
\end{corollary}
\noindent The mechanism is a collision: at every positive multiple of $L$ at
least $|F|$ support divisors land on the same index, so $f_A\ge|F|$ there,
while the exact carry equation $f_A(n)+e_n=k_{n-1}+2e_{n-1}$ with
$k_{n-1}\le1$ lets one step absorb at most one unit --- so the preceding
excess must spike to at least $\lceil(|F|-1)/2\rceil$, and the spikes recur
with density $1/L$. Formalised:
\lword{RationalSupportCarrySkeleton.lean}{2037}{booleanCollision_wrap_bound_of_common_multiple}{common-multiple wrap bound}
(denominator-shifted form
\lloc{RationalSupportCarrySkeleton.lean}{1863}),
the carry equation
\lword{RationalSupportCarrySkeleton.lean}{1651}{supportCoeff_add_oddTailExcess_succ_eq_wrap_add_two_excess}{wrap--excess recurrence},
and the dyadic corollary
\lword{RationalSupportCarrySkeleton.lean}{2200}{dyadic_support_fraction_reciprocalMass_diverges_or_gt_one}{dyadic reciprocal-mass dichotomy}
(via \lloc{RationalSupportCarrySkeleton.lean}{2003}).

\begin{theorem}[carry states are unbounded over infinite supports]\label{res:unbounded}
Let $A$ be infinite with value $p/(2^c v)$, $v\ge1$. Then the positive integer
carry state $u(N)=v\,T_{f_A}(c+N)$ is unbounded: for every $B$ there is an $N$
with $u(N)>B$.
\end{theorem}
\noindent The proof picks $2B+1$ positive support elements; at a common
multiple $L$ of all of them, the local inequality $1+v\,|F|\le2\,u(L-c-1)$
pushes the state past $B$. Formalised:
\lword{RationalSupportCarrySkeleton.lean}{2377}{exists_unbounded_shifted_odd_tail_nat_state_of_support_fraction}{unbounded shifted states},
from
\lword{RationalSupportCarrySkeleton.lean}{2320}{shifted_state_unbounded_of_infinite_support}{infinite-support state unboundedness}
and the local inequality
\lword{RationalSupportCarrySkeleton.lean}{2228}{one_add_mul_card_le_two_mul_shifted_state}{support-cardinality state bound};
the $\{2,3\}$ example revalidates it at sixty-six common multiples
(\lword{RationalSupportCarrySkeleton.lean}{2401}{carryOrbit23_common_multiple_bound_validation66}{two-point finite validation}).
So a rational value over an infinite support cannot run on a bounded carry
alphabet. That closes finite-state readings of the carry system; it does not
close the problem. Nothing in the constraints above
prevents an infinite support from satisfying every one of these constraints.
The repetend identities and divisor averages used here are classical. We make
no claim of mathematical novelty for the coupled reciprocal-mass bounds,
collision strengthening, or global unboundedness statement.

\subsection{Greedy geometry of the \#257 value set}

Finally, consider the set of all values a base-$2$ support series can
take. For $n\ge1$ put $w_n=1/(2^n-1)$, let $T(n)$ be the mass strictly after
exponent $n$, and let
\[
\mathcal A\;=\;\Bigl\{\,\sum_{n\in A}w_n \;:\; A\subseteq\N,\ 0\notin A\Bigr\},
\]
the \emph{Mersenne achievement set}, the exponent $0$ being normalised away as
analytically invisible. Coding by supports of positive exponents is literally
the formal support series
(\lword{GreedyAchievementSet.lean}{554}{positiveMersenneSupportValue_eq_erdosSupportSeries}{support-value identity};
the set is \lword{GreedyAchievementSet.lean}{571}{mersenneAchievementSet}{achievement-set definition}). In
this language the base-$2$ \#257 statement reads: the rational members of
$\mathcal A$ are exactly the finite subset sums. Nothing below decides that.
Kova\v{c} and Tao verify, for every fixed integer $t\ge2$, that
\[
  \sum_{\ell>n}\frac1{t^\ell-1}<\frac1{t^n-1},
\]
and deduce distinct Lambert subsums and a Cantor set~\cite{kovactao}. At
$t=2$ this is the relevant Lambert-series instance of Kakeya's classical
strict-tail criterion for subsum sets~\cite{kakeya1914}. The local
contribution is the checked greedy criterion, quantitative enclosure, and
finite non-membership certificates over that known geometry; it does not
attribute those refinements to the cited sources or settle \#257.

\begin{theorem}[superincreasing geometry and greedy membership]\label{res:greedy-details}
For every $n\ge1$ the weight dominates the whole tail after it, $T(n)<w_n$,
with the two-scale gap $w_n-T(n)=\tfrac23\,4^{-n}+O(8^{-n})$ and explicit
valid $O$-constant $3$. Consequently a real $x$ belongs to $\mathcal A$
exactly when $x\ge0$ and every greedy residual of $x$ is at most the remaining
tail; and normalised support coding is injective, so each member of
$\mathcal A$ has a unique support, which the greedy recursion recovers bit by
bit.
\end{theorem}
\noindent Formalised:
\lword{GreedyAchievementSet.lean}{180}{mersenneTail_lt_weight}{strict-tail inequality},
\lword{GreedyAchievementSet.lean}{473}{mersenneGap_isBigOWith}{gap big-O estimate} (explicit bound
\lword{GreedyAchievementSet.lean}{448}{mersenneGap_asymptotic_bound}{explicit gap bound}),
\lword{GreedyAchievementSet.lean}{1445}{mem_mersenneAchievementSet_iff_greedy_survival}{greedy survival criterion},
and
\lword{GreedyAchievementSet.lean}{1527}{positiveMersenneSupportValue_injective_normalized}{normalised support-value injectivity},
transferred to the kernel series at
\lword{GreedyAchievementSet.lean}{1536}{erdosSupportSeries_two_injective_normalized}{normalised support-series injectivity}.
The nonnegativity guard in the membership criterion is necessary: without it a
negative target would survive every level while lying outside $\mathcal A$.

The binary coding also permits the global geometry to be checked directly.

\begin{proposition}[fat-Cantor geometry]\label{res:greedytopology}
The achievement set $\mathcal A$ is compact, perfect, totally disconnected,
and nowhere dense.  Its Lebesgue measure is exactly one.
\end{proposition}
\noindent Formalised:
\lword{GreedyAchievementSet.lean}{996}{volume_mersenneAchievementSet}{measure-one theorem},
\lword{GreedyAchievementSet.lean}{1620}{perfect_mersenneAchievementSet}{perfectness},
\lword{GreedyAchievementSet.lean}{1636}{isTotallyDisconnected_mersenneAchievementSet}{total disconnectedness},
and
\lword{GreedyAchievementSet.lean}{1645}{isNowhereDense_mersenneAchievementSet}{nowhere density}.
Kova\v{c}--Tao supply the strict-tail Cantor-set context cited above.  The
formalisation records these exact properties without making a novelty or
priority claim for them.

\begin{proposition}[exact rational runs and finite non-membership certificates]\label{res:death}
The greedy recursion executed in exact rational arithmetic agrees step by step
with the real recursion under casting. A finite decidable certificate --- the
exact rational residual at some level is at least an exact rational upper
enclosure of the remaining tail --- soundly proves non-membership in
$\mathcal A$; the level-one instance shows $3/4\notin\mathcal A$. And a
rational number already known to lie in $\mathcal A$ has computable support
bits.
\end{proposition}
\noindent Formalised:
\lword{GreedyAchievementSet.lean}{1130}{cast_greedyMersenneRemainderRat}{rational--real greedy agreement};
\lword{GreedyAchievementSet.lean}{1726}{certifiedGreedyMersenneDeath_not_mem}{finite death certificate},
with the certificate
\lword{GreedyAchievementSet.lean}{1720}{CertifiedGreedyMersenneDeath}{death-certificate predicate} and the
enclosure \lword{GreedyAchievementSet.lean}{1687}{mersenneTailUpperRat}{rational tail enclosure};
\lword{GreedyAchievementSet.lean}{1748}{three_fourths_not_mem_mersenneAchievementSet}{three-quarters nonmembership};
and \lword{GreedyAchievementSet.lean}{1832}{rational_member_support_bit_iff}{rational-member support-bit equivalence}.
The certificates are one-sided: survival through any finite depth proves
nothing about membership, and no decidability of membership in $\mathcal A$ is
asserted.  Neither the global geometry nor the finite death certificates
exclude rational values for either open series.

\paragraph{Appendix consequence.}
Appendices~B.1--B.2 give exact normal forms for rationality;
Appendices~B.3--B.4 impose necessary asymptotic constraints on any infinite
rational support; and Appendix~B.5 gives the canonical achievement-set
representation and finite nonmembership certificates.  None of these
necessary conditions is contradictory.  Their role is to make the rational
side exact and auditable, not to close either open problem.

\section{Technical reductions for the totient-certificate branch}
\label{app:totient-technical}

This appendix expands the supporting mathematics behind
Theorem~\ref{res:red21}, Proposition~\ref{res:complete}, and
Corollary~\ref{res:diag}.  It contains exact cone-flatness and target-interval
reformulations, finite residue projections, denominator and bit-lifting
formulae, bounded certificate families, and rigorous obstructions to several
candidate strategies.  Each statement retains its own hypotheses and status.
None proves the unbounded certificate supply
\eqref{eq:unbounded-certificate-supply}.

The dependency order is: exact normal forms in Appendix~C.1; finite
projections and verified instances in Appendix~C.2; fresh-prime and bit-level
requirements, including scoped no-go results, in Appendix~C.3; and finite
determinant and rank routes in Appendix~C.4.  The first two say exactly what
can be checked at one scale; the latter two delimit what would be needed to
make those checks unbounded.

\TotientTechnicalAppendix

\section{Lambert, probability, and finite-algebraic complements}
\label{app:lambert-complements}

This appendix records results that enlarge the surrounding Lambert-series
picture without advancing either unbounded frontier in
Section~\ref{sec:further}.  They include periodic-weight theorems, finite
denominator obstructions, the M\"obius-square identity, a squared Lambert
calculus for gcd moments, an exact Stern--Brocot cylinder law, and sharp
Fibonacci--continuant stability for alternating runs.

\FurtherLambertAppendix

\paragraph{Appendix consequence.}
This appendix contains three different statuses.  The periodic-weight
irrationality theorems are unconditional but lie within the cited classical
and modern theory; the M\"obius-square, gcd-moment, cylinder, and continuant
statements are unconditional structural results; and the finite-algebraic
propositions delimit strategies without producing an irrationality
criterion.  None supplies the unbounded certificate or half-membership input
isolated in Section~\ref{sec:further}.

\section{Guide to the formal sources}\label{app:index}

The paper is meant to be read as mathematics. The blue mathematical phrases
following formalised statements are optional links to the exact checked
theorems in the \href{\sourceurl}{archived source revision}. Exact declaration
names remain in the machine-readable audit layer.

For a first reading, the main landmarks are Theorem~\ref{res:full} for the
full-support irrationality theorem, Theorem~\ref{res:farey} for the totient
denominator exclusion, Theorem~\ref{res:red21} for the exact
tail-certificate reduction, and Proposition~\ref{res:greedytopology} for the
geometry of the Mersenne achievement set.  For the strongest independent
structure beyond those two proof spines, see
Theorems~\ref{res:sternbrocotcylinders} and~\ref{res:sternbrocotruns}.
The public repository contains a complete searchable cross-reference for
readers who need declaration-level detail.

% Retain the exhaustive source-link manifest in the authored TeX for automated
% validation and query tooling, but keep it outside the printed paper.
\iffalse

Each row links to the declaration at the archived source revision.

\begin{center}
\footnotesize
\begin{tabular}{@{}p{7.2cm}p{7.8cm}@{}}
\toprule
informal statement & linked Lean source\\
\midrule
$\sum 1/(b^n-1)$ irrational, all $b\ge2$ & \lrefx{CertificateKernel.lean}{8000}{irrational_erdosSum_full_support}\ (\texttt{CertificateKernel})\\
Erd\H{o}s--Borwein $E$ irrational & \lrefx{CertificateKernel.lean}{8007}{irrational_erdosBorwein_series}\ (\texttt{CertificateKernel})\\
factorial support & \lrefx{CertificateKernel.lean}{5707}{irrational_erdosSum_factorial_support}\ (\texttt{CertificateKernel})\\
power-of-two support & \lrefx{CertificateKernel.lean}{5731}{irrational_erdosSum_two_pow_support}\ (\texttt{CertificateKernel})\\
pairwise-coprime supports & \lrefx{CertificateKernel.lean}{10448}{irrational_erdosSupportSeries_pairwise_coprime}\ (\texttt{CertificateKernel})\\
eventually-periodic nonnegative rational weights & \lrefx{CertificateKernel.lean}{12483}{irrational_ratWeightSeries_eventuallyPeriodic}\ (\texttt{CertificateKernel})\\
periodic signed-weight irrational-or-terminating dichotomy & \lrefx{CertificateKernel.lean}{13847}{irrational_or_bpow_mul_eq_intCast_intWeightedErdosSeries_periodic}\ (\texttt{CertificateKernel})\\
period-four residue obstruction & \lrefx{CertificateKernel.lean}{13898}{intWeightedCoeff_periodFourSignWeight_eq_zero_of_mod_four_eq_three}\ (\texttt{CertificateKernel})\\
$L(\mu)=1/2$ & \lrefx{MersenneLambertLadder.lean}{587}{tsum_moebius_div_two_pow_sub_one_eq_half}\ (\texttt{MersenneLambertLadder})\\
$L(\ph)=2$ & \lrefx{MersenneLambertLadder.lean}{575}{tsum_totient_div_two_pow_sub_one_eq_two}\ (\texttt{MersenneLambertLadder})\\
positive lift $L(A)=S$ & \lrefx{CertificateKernel.lean}{18117}{tsum_primWeight_div_two_pow_sub_one_eq_totient_series}\ (\texttt{CertificateKernel})\\
M\"obius-square lens & \lrefx{CertificateKernel.lean}{18126}{totient_series_eq_half_add_moebius_mersenne_square}\ (\texttt{CertificateKernel})\\
$L_2(\mathbf1)$ and $L_2(\ph)$ gcd moments & \lrefx{GcdMomentCalculus.lean}{216}{tsum_one_div_mersenne_sq_eq_sigma_sub_tau_series}, \lrefx{GcdMomentCalculus.lean}{235}{tsum_totient_div_mersenne_sq_eq_gcd_moment_series}\ (\texttt{GcdMomentCalculus})\\
Stern--Brocot cylinder recursion and rate & \lrefx{GcdMomentCalculus.lean}{474}{cylinderMass_split}, \lrefx{GcdMomentCalculus.lean}{525}{sternBrocotDepthMass_error}\ (\texttt{GcdMomentCalculus})\\
$S=\tfrac12+\Pr(\gcd=1)$ & \lrefx{CertificateKernel.lean}{18229}{totient_series_eq_half_add_visible_coprime_pairs}\ (\texttt{CertificateKernel})\\
$S\ne p/q$ for $q\le\Qzero$ & \lrefx{CertificateKernel.lean}{18056}{tsum_totient_div_pow_two_ne_ratCast_of_den_le_79639646646701375323355774875831053}\ (\texttt{CertificateKernel})\\
tail-period reduction & \lrefx{TotientTailPeriodKiller.lean}{394}{irrational_totient_series_of_certificate_supply}\ (\texttt{TotientTailPeriodKiller})\\
diagonal reduction & \lrefx{LcmDiagonalReduction.lean}{112}{irrational_totient_series_of_lcm_diagonal_certificate_supply}\ (\texttt{LcmDiagonalReduction})\\
cone flatness from rationality & \lrefx{LcmConeFlatness.lean}{90}{rational_totient_series_forces_lcm_cone_flatness}\ (\texttt{LcmConeFlatness})\\
certificate completeness & \lrefx{LcmConeFlatness.lean}{316}{exists_certifiedKill_iff_tail_diff_notMem_int}\ (\texttt{LcmConeFlatness})\\
joint cone refuter & \lrefx{LcmConeNonflat.lean}{126}{exists_nonintegral_pair_of_coneNonflatCert}\ (\texttt{LcmConeNonflat})\\
verified finite instances & \lrefx{TotientTailPeriodKiller.lean}{404}{certifiedKill_all_small}\ (\texttt{TotientTailPeriodKiller})\\
full-target diagonal reduction & \lrefx{DiagonalPincerDecomposition.lean}{215}{diagonal_int_iff_foreignDiagonalDefect_hits_fullTarget}\ (\texttt{DiagonalPincerDecomposition})\\
squared-Mersenne enclosure & \lrefx{SquaredMersenneDiagonalEnclosure.lean}{116}{scaleDiagonalTailDifference_sub_lambertProjectedDiagonal}\ (\texttt{SquaredMersenneDiagonalEnclosure})\\
exact Mersenne-shadow denominator & \lrefx{MersenneShadowDenominatorGrowth.lean}{147}{lcmHeight_scaledMobiusShadow_den_exact}\ (\texttt{MersenneShadowDenominatorGrowth})\\
fresh-loss projection reduction & \lrefx{DiagonalFreshLossBridge.lean}{191}{irrational_totientSeries_of_diagonalFreshLossProjectionSupply}\ (\texttt{DiagonalFreshLossBridge})\\
prime-adjunction diamond no-go & \lrefx{FullTargetPrimeAdjunctionNoGo.lean}{196}{fullTarget_primeAdjunction_diamond_iff_root}\ (\texttt{FullTargetPrimeAdjunctionNoGo})\\
scalar-localisation height obstruction & \lrefx{AdelicHeightObstruction.lean}{23}{scalarLocalization_complement_dvd}\ (\texttt{AdelicHeightObstruction})\\
square-CRT correction suppression & \lrefx{SquareCRTCube.lean}{198}{exists_squareCRT_clean_totient_family}\ (\texttt{SquareCRTCube})\\
signed dyadic-moment algebra & \lrefx{SignedQMomentObstruction.lean}{28}{det_mul_rectangular}\ (\texttt{SignedQMomentObstruction})\\
dyadic totient certificate rank criterion & \lrefx{TotientMahlerDefect.lean}{98}{finrank_canonicalTotientKernel_eq_of_certificate}\ (\texttt{TotientMahlerDefect})\\
residual-gauge obstruction for monomial minors & \lrefx{ResidualGaugeObstruction.lean}{55}{det_residualMonomialMatrix_ne_zero_iff}\ (\texttt{ResidualGaugeObstruction})\\
tail-orbit rationality criterion & \lrefx{GenericTailOrbitRigidity.lean}{425}{binaryCoeffSeries_rational_iff_exists_temperedBinaryOrbit}\ (\texttt{GenericTailOrbitRigidity})\\
Boolean carry certificate $\Leftrightarrow$ support fraction & \lrefx{BooleanMobiusCarry.lean}{948}{exists_normalized_support_fraction_iff_exists_booleanMobiusCarry}\ (\texttt{BooleanMobiusCarry})\\
one-wrap classification & \lrefx{RationalSupportCarrySkeleton.lean}{415}{one_doublingWrap_classification}\ (\texttt{RationalSupportCarrySkeleton})\\
order-wrap mass bound & \lrefx{RationalSupportCarrySkeleton.lean}{1478}{one_div_oddOrder_le_reciprocalMass_of_support_fraction}\ (\texttt{RationalSupportCarrySkeleton})\\
unbounded carry states & \lrefx{RationalSupportCarrySkeleton.lean}{2377}{exists_unbounded_shifted_odd_tail_nat_state_of_support_fraction}\ (\texttt{RationalSupportCarrySkeleton})\\
sublogarithmic zero windows & \lrefx{SublogDivisorCoverage.lean}{392}{supportCoeffZeroWindow_length_le_eps_logb_add}\ (\texttt{SublogDivisorCoverage})\\
greedy membership criterion & \lrefx{GreedyAchievementSet.lean}{1445}{mem_mersenneAchievementSet_iff_greedy_survival}\ (\texttt{GreedyAchievementSet})\\
Mersenne achievement set: measure one and nowhere dense & \lrefx{GreedyAchievementSet.lean}{996}{volume_mersenneAchievementSet}, \lrefx{GreedyAchievementSet.lean}{1645}{isNowhereDense_mersenneAchievementSet}\ (\texttt{GreedyAchievementSet})\\
Campbell shift-synchronization boundary & \lrefx{CampbellShiftSynchronization.lean}{299}{campbellQuarterExponent_periodFreeze_noSync}, \lrefx{CampbellShiftSynchronization.lean}{488}{shiftWindowZero_iff_half_mem_mersenneAchievementSet}\ (\texttt{CampbellShiftSynchronization})\\
\bottomrule
\end{tabular}
\end{center}
\fi

\clearpage
\begin{thebibliography}{99}
\small
\setlength{\itemsep}{0pt}
\bibitem{erdos1948} P.~Erd\H{o}s, \emph{On arithmetical properties of Lambert series},
  J. Indian Math. Soc. (N.S.) \textbf{12} (1948), 63--66,
  \href{https://users.renyi.hu/~p_erdos/1948-04.pdf}{archive scan}.
\bibitem{erdos1968} P.~Erd\H{o}s, \emph{On the irrationality of certain series},
  Math. Student \textbf{36} (1968), 222--226,
  \href{https://users.renyi.hu/~p_erdos/1969-09.pdf}{archive scan}.
\bibitem{erdosgraham1980} P.~Erd\H{o}s and R.~L. Graham, \emph{Old and New Problems and
  Results in Combinatorial Number Theory}, Monographies de l'Enseignement
  Math\'ematique, vol.~28, Universit\'e de Gen\`eve, Geneva, 1980, pp.~61--62,
  \href{https://mathweb.ucsd.edu/~ronspubs/80_11_number_theory.pdf}{author-hosted scan}.
\bibitem{erdos1988} P.~Erd\H{o}s, \emph{On the irrationality of certain series: problems and
  results}, in \emph{New Advances in Transcendence Theory}, ed.~A.~Baker,
  Cambridge University Press, 1988, pp.~102--109,
  doi:\href{https://doi.org/10.1017/CBO9780511897184.009}{10.1017/CBO9780511897184.009}.
\bibitem{borwein1992} P.~B. Borwein, \emph{On the irrationality of certain series},
  Math. Proc. Cambridge Philos. Soc. \textbf{112} (1992), no.~1, 141--146,
  doi:\href{https://doi.org/10.1017/S030500410007081X}{10.1017/S030500410007081X}.
\bibitem{apostol} T.~M. Apostol, \emph{Introduction to Analytic Number Theory},
  Undergraduate Texts in Mathematics, Springer, 1976,
  doi:\href{https://doi.org/10.1007/978-1-4757-5579-4}{10.1007/978-1-4757-5579-4}.
\bibitem{merca2017} M.~Merca, \emph{The Lambert series factorization theorem},
  Ramanujan J. \textbf{44} (2017), no.~2, 417--435,
  doi:\href{https://doi.org/10.1007/s11139-016-9856-3}{10.1007/s11139-016-9856-3}.
\bibitem{mercaschmidt} M.~Merca and M.~D. Schmidt, \emph{Generating special arithmetic
  functions by Lambert series factorizations}, Contrib. Discrete Math. \textbf{14}
  (2019), no.~1, 31--45,
  doi:\href{https://doi.org/10.55016/ojs/cdm.v14i1.62425}{10.55016/ojs/cdm.v14i1.62425}.
\bibitem{kakeya1914} S.~Kakeya, \emph{On the set of partial sums of an infinite series},
  Proc. Tokyo Math.-Phys. Soc., 2nd ser. \textbf{7} (1914), 250--251,
  doi:\href{https://doi.org/10.11429/ptmps1907.7.14_250}{10.11429/ptmps1907.7.14\_250}.
\bibitem{farey1816} J.~Farey, \emph{On a curious property of vulgar fractions},
  Philos. Mag. \textbf{47} (1816), 385--386,
  doi:\href{https://doi.org/10.1080/14786441608628487}{10.1080/14786441608628487}.
\bibitem{nesterenko} Yu.~V. Nesterenko, \emph{Modular functions and transcendence questions},
  Mat. Sb. \textbf{187} (1996), no.~9, 65--96; English transl., Sb. Math. \textbf{187} (1996), no.~9, 1319--1348,
  doi:\href{https://doi.org/10.1070/SM1996v187n09ABEH000158}{10.1070/SM1996v187n09ABEH000158}.
\bibitem{postelmansvanassche} K.~Postelmans and W.~Van Assche,
  \emph{Irrationality of $\zeta_q(1)$ and $\zeta_q(2)$},
  J. Number Theory \textbf{126} (2007), no.~1, 119--154,
  doi:\href{https://doi.org/10.1016/j.jnt.2006.11.011}{10.1016/j.jnt.2006.11.011}.
\bibitem{lucatachiya} F.~Luca and Y.~Tachiya, \emph{Irrationality of Lambert series associated
  with a periodic sequence}, Int. J. Number Theory \textbf{10} (2014), no.~3, 623--636,
  doi:\href{https://doi.org/10.1142/S1793042113501121}{10.1142/S1793042113501121}.
\bibitem{kovactao} V.~Kova\v{c} and T.~Tao, \emph{On several irrationality problems for
  Ahmes series}, Acta Math. Hungar. \textbf{175} (2025), 572--608,
  doi:\href{https://doi.org/10.1007/s10474-025-01528-0}{10.1007/s10474-025-01528-0}.
\bibitem{wanggrauribas} H.~Wang and J.~M. Grau Ribas, \emph{Positive dyadic density for
  rational weighted binary expansions}, arXiv:2606.24972 (2026),
  doi:\href{https://doi.org/10.48550/arXiv.2606.24972}{10.48550/arXiv.2606.24972}.
\bibitem{crandall} R.~Crandall, \emph{The googol-th bit of the Erd\H{o}s--Borwein constant},
  Integers \textbf{12} (2012), no.~5, 811--840,
  doi:\href{https://doi.org/10.1515/integers-2012-0007}{10.1515/integers-2012-0007}.
\bibitem{campbell} J.~M. Campbell, \emph{On the binary digits of the Erd\H{o}s--Borwein
  constant}, arXiv:2605.24160 (2026),
  doi:\href{https://doi.org/10.48550/arXiv.2605.24160}{10.48550/arXiv.2605.24160}.
\bibitem{taoteravainen} T.~Tao and J.~Ter\"av\"ainen, \emph{Quantitative correlations and
  some problems on prime factors of consecutive integers}, arXiv:2512.01739v2 (2026),
  doi:\href{https://doi.org/10.48550/arXiv.2512.01739}{10.48550/arXiv.2512.01739}.
\bibitem{lean} L.~de Moura, S.~Kong, J.~Avigad, F.~van Doorn, and J.~von Raumer,
  \emph{The Lean theorem prover (system description)}, in \emph{CADE-25}, LNCS 9195,
  Springer (2015), 378--388, doi:10.1007/978-3-319-21401-6\_26.
\bibitem{mathlib} The mathlib Community, \emph{The Lean mathematical library},
  in \emph{CPP 2020}, ACM (2020), 367--381, doi:10.1145/3372885.3373824.
\bibitem{erdosproblems} T.~Bloom (ed.), \emph{Erd\H{o}s Problems}, entries 249 and 257,
  \href{https://www.erdosproblems.com}{\texttt{erdosproblems.com}}, accessed July 2026.

\end{thebibliography}

\end{document}
